\documentclass[11pt,a4paper]{article}

% ─── Packages ─────────────────────────────────────────────────────────────────
\usepackage[T1]{fontenc}
\usepackage[utf8]{inputenc}
\usepackage{lmodern}
\usepackage[margin=2.4cm]{geometry}
\usepackage{hyperref}
\usepackage{xcolor}
\usepackage{listings}
\usepackage{booktabs}
\usepackage{longtable}
\usepackage{enumitem}
\usepackage{multicol}
\usepackage{fancyhdr}
\usepackage{titlesec}
\usepackage{xspace}

\hypersetup{
  colorlinks=true,
  linkcolor=blue!60!black,
  urlcolor=blue!60!black,
  citecolor=blue!60!black,
}

% ─── Code listing styles ─────────────────────────────────────────────────────
\definecolor{codebg}{gray}{0.96}
\definecolor{coqkw}{rgb}{0.0,0.4,0.0}
\definecolor{ckw}{rgb}{0.0,0.0,0.6}
\definecolor{cmt}{rgb}{0.5,0.5,0.5}
\definecolor{str}{rgb}{0.6,0.0,0.0}

\lstdefinestyle{rocqstyle}{
  language={[Objective]Caml},
  backgroundcolor=\color{codebg},
  basicstyle=\ttfamily\small,
  keywordstyle=\color{coqkw}\bfseries,
  commentstyle=\color{cmt}\itshape,
  stringstyle=\color{str},
  numbers=left,
  numberstyle=\tiny\color{gray},
  numbersep=6pt,
  xleftmargin=2em,
  frame=l,
  framesep=4pt,
  framerule=1pt,
  rulecolor=\color{coqkw!40},
  breaklines=true,
  breakatwhitespace=true,
  showstringspaces=false,
  tabsize=2,
  morekeywords={Definition,Fixpoint,match,with,end,let,in,if,then,else,
    fun,forall,Inductive,Record,Type,Set,Prop,Some,None,Step,Error,Halt,
    CCall_request},
  literate={=>}{$\Rightarrow$}2 {:=}{$\coloneq$}2,
}

\lstdefinestyle{cstyle}{
  language=C,
  backgroundcolor=\color{codebg},
  basicstyle=\ttfamily\small,
  keywordstyle=\color{ckw}\bfseries,
  commentstyle=\color{cmt}\itshape,
  stringstyle=\color{str},
  numbers=left,
  numberstyle=\tiny\color{gray},
  numbersep=6pt,
  xleftmargin=2em,
  frame=l,
  framesep=4pt,
  framerule=1pt,
  rulecolor=\color{ckw!40},
  breaklines=true,
  breakatwhitespace=true,
  showstringspaces=false,
  tabsize=2,
  morekeywords={value,intnat,uintnat,mlsize_t,code_t,tag_t,opcode_t,
    Instruct,Next,Field,Val_int,Val_long,Long_val,Code_val,Closinfo_val,
    Make_closinfo,Make_header,Alloc_small,Atom,Is_block,Tag_val,Wosize_val,
    Double_wosize,Double_array_tag,Closure_tag,Infix_tag,
    Setup_for_c_call,Restore_after_c_call,
    Setup_for_gc,Restore_after_gc,Val_unit,Val_false,Val_not,
    Byte_u,Int_val,Store_double_flat_field,Double_val,Double_flat_field,
    Store_double_val,Trap_pc,Trap_link_offset,
    Caml_state,caml_modify,caml_initialize},
}

\lstnewenvironment{rocqcode}[1][]{\lstset{style=rocqstyle,#1}}{}
\lstnewenvironment{ccode}[1][]{\lstset{style=cstyle,#1}}{}

% ─── Custom commands ─────────────────────────────────────────────────────────

% Instruction entry heading
\newcommand{\instruction}[2]{%
  \subsection{\texttt{#1}}%
  \label{instr:#2}%
}

% Status badges
\newcommand{\statusok}{\textbf{\textcolor{green!50!black}{OK}}\xspace}
\newcommand{\statusbug}[1]{\textbf{\textcolor{red!70!black}{BUG:}} #1\xspace}
\newcommand{\statusneedstest}[1]{\textbf{\textcolor{orange!80!black}{NEEDS TEST:}} #1\xspace}
\newcommand{\statusstub}{\textbf{\textcolor{gray}{STUB}}\xspace}

% Risk badges
\newcommand{\risklow}{\textcolor{green!50!black}{\textbf{Low}}\xspace}
\newcommand{\riskmed}{\textcolor{orange!80!black}{\textbf{Medium}}\xspace}
\newcommand{\riskhigh}[1]{\textcolor{red!70!black}{\textbf{High:}} #1\xspace}

% Safe input: include file if it exists, else show placeholder
\newcommand{\inputsnippet}[1]{%
  \IfFileExists{#1}{\input{#1}}{%
    \begin{quote}\itshape [Snippet not yet extracted: \texttt{#1}]\end{quote}%
  }%
}

% Short aliases
\newcommand{\rocqsnippet}[1]{\inputsnippet{snippets/rocq/#1.tex}}
\newcommand{\csnippet}[1]{\inputsnippet{snippets/c/#1.tex}}

% ─── Document ─────────────────────────────────────────────────────────────────
\title{\textbf{Audit: Rocq Bytecode Interpreter vs.\ OCaml 4.14 \texttt{interp.c}}\\[0.3em]
  \large Systematic Per-Instruction Comparison}
\author{Generated by audit infrastructure}
\date{\today}

\begin{document}
\maketitle
\tableofcontents
\newpage

%═══════════════════════════════════════════════════════════════════════════════
\section{Introduction}
%═══════════════════════════════════════════════════════════════════════════════

\subsection{Purpose}

This document provides a systematic, per-instruction audit of the formally
verified OCaml bytecode interpreter written in Rocq (\texttt{Interp.v}) against
two authoritative references:

\begin{enumerate}
  \item \textbf{OCaml 4.14.2 \texttt{runtime/interp.c}} --- the production C
    implementation of the ZINC bytecode interpreter.
  \item \textbf{The Cadmium bytecode instruction specification}
    (\texttt{caml-instructions.pdf}) --- a detailed reference manual for
    each instruction's semantics.
\end{enumerate}

The interpreter is the core \emph{trusted} component of the verified compiler
project.  Any discrepancy between it and the real OCaml runtime constitutes a
soundness risk for the entire verification effort.

\subsection{Methodology}

For each of the 107 instruction variants defined in \texttt{AST.v}, we:

\begin{enumerate}
  \item Extract the exact C code from \texttt{interp.c} (OCaml 4.14.2).
  \item Extract the exact Rocq code from \texttt{Interp.v}.
  \item Paraphrase the PDF specification.
  \item Compare the three, noting any discrepancies.
  \item Assign a \emph{status} (\statusok, bug, or needs-test) and a
    \emph{risk level} (low, medium, or high).
\end{enumerate}

Code snippets are auto-extracted by \texttt{extract\_snippets.py} so
the document stays current as \texttt{Interp.v} evolves.

\subsection{Scope}

\begin{itemize}
  \item \textbf{In scope:} All 107 instruction variants in \texttt{AST.v},
    the machine state model (\texttt{Machine.v}), the value representation
    (\texttt{Value.v}), and helper functions (\texttt{do\_raise},
    \texttt{z\_flip\_sign}, \texttt{field\_or\_heap}, etc.).
  \item \textbf{Out of scope:} GC interactions, threading, signal handling
    (modeled as no-ops), native code compilation, and debugger integration.
\end{itemize}

\subsection{Key Architectural Differences}

\begin{longtable}{p{3.5cm}p{5.5cm}p{5.5cm}}
\toprule
\textbf{Aspect} & \textbf{OCaml \texttt{interp.c}} & \textbf{Rocq \texttt{Interp.v}} \\
\midrule
\endhead
Stack & C pointer \texttt{sp} (grows down) & Functional list (grows with cons) \\
Heap & GC-managed C heap & \texttt{PositiveMap} (immutable map) \\
Values & Tagged words (tag bit) & Algebraic type (\texttt{Val\_int}, \texttt{Val\_block}, etc.) \\
Closures & Pointer + infix offset via pointer arithmetic & \texttt{Val\_closure addr ofs} \\
PC & C \texttt{code\_t} pointer & \texttt{Z} index into \texttt{PrimArray} \\
Trap pointer & C pointer \texttt{Caml\_state->trapsp} & \texttt{trap\_sp : nat} = stack depth \\
C calls & Direct function pointer call & \texttt{CCall\_request} step result \\
GC & Triggered by allocation & No GC (immutable heap) \\
Signals & \texttt{caml\_something\_to\_do} & Ignored (no-op) \\
Word size & 63-bit integers (tagged) & Arbitrary-precision \texttt{Z} (extraction uses 63-bit) \\
\bottomrule
\end{longtable}


%═══════════════════════════════════════════════════════════════════════════════
\section{Machine Model}
\label{sec:machine}
%═══════════════════════════════════════════════════════════════════════════════

\subsection{State Record}

The Rocq machine state (\texttt{Machine.v}):

\begin{rocqcode}
Record state : Type := mk_state {
  pc         : Z;
  accu       : value;
  stack      : list value;
  env        : value;
  extra_args : nat;
  global     : list value;
  trap_sp    : nat;
  hp         : heap;
  next_addr  : nat;
}.
\end{rocqcode}

This corresponds to the C registers:
\begin{ccode}
register code_t pc;
register value * sp;      // our stack : list value
register value accu;
value env;
intnat extra_args;
// Caml_state->trapsp      // our trap_sp : nat
// caml_global_data         // our global : list value
\end{ccode}

\paragraph{Analysis.}
The correspondence is direct.  The main difference is that the C \texttt{sp}
is a downward-growing pointer while our \texttt{stack} is an
upward-growing cons-list.  Both represent the same abstract stack.
The \texttt{trap\_sp} field stores the stack \emph{depth} at the time of the
last \texttt{PUSHTRAP}, while \texttt{interp.c} stores an absolute pointer.
The Rocq version must compute the offset via
\texttt{length(stack) - trap\_sp} to find the trap frame, which is equivalent.

\subsection{Heap Model}

The heap is a \texttt{PositiveMap.t (nat * list value)}, mapping addresses to
(tag, fields) pairs.  Allocation returns monotonically increasing addresses.
There is no garbage collection; the heap only grows.

\subsection{Value Representation}

\begin{rocqcode}
Inductive value : Type :=
  | Val_int : Z -> value
  | Val_block : nat -> list value -> value
  | Val_ptr : nat -> value
  | Val_closure : nat -> nat -> value.
\end{rocqcode}

In \texttt{interp.c}, the tag bit distinguishes integers from pointers.
Our algebraic representation is unambiguous by construction.  The
\texttt{Val\_closure} constructor explicitly tracks the base address and infix
offset within a shared closure block, corresponding to the C pointer arithmetic
\texttt{env + offset * sizeof(value)}.


%═══════════════════════════════════════════════════════════════════════════════
\section{Stack Manipulation}
\label{sec:stack}
%═══════════════════════════════════════════════════════════════════════════════

% ─── ACC ──────────────────────────────────────────────────────────────────────
\instruction{ACC}{acc}

\paragraph{PDF Spec.}
\texttt{ACC $n$}: Sets \texttt{accu} to \texttt{sp[$n$]}.  The specialized
variants \texttt{ACC0}--\texttt{ACC7} are optimized forms.

\paragraph{OCaml \texttt{interp.c} (line 322).}
\csnippet{ACC}

\paragraph{Rocq \texttt{Interp.v} (line 110).}
\rocqsnippet{ACC}

\paragraph{Analysis.}
Direct correspondence.  The C \texttt{sp[$n$]} maps to
\texttt{nth\_error stack $n$}.  Our version returns \texttt{Error} on
out-of-bounds; the C version has undefined behavior (buffer overread).

\paragraph{Status.} \statusok

\paragraph{Risk.} \risklow


% ─── PUSH ─────────────────────────────────────────────────────────────────────
\instruction{PUSH}{push}

\paragraph{PDF Spec.}
\texttt{PUSH}: Pushes \texttt{accu} onto the stack (\texttt{*-{}-sp = accu}).
Same as \texttt{PUSHACC0}.

\paragraph{OCaml \texttt{interp.c} (line 339).}
\csnippet{PUSH}

\paragraph{Rocq \texttt{Interp.v} (line 116).}
\rocqsnippet{PUSH}

\paragraph{Analysis.}
\texttt{PUSH} conses \texttt{accu} onto the stack.  Exact match.

\paragraph{Status.} \statusok

\paragraph{Risk.} \risklow


% ─── PUSHACC ──────────────────────────────────────────────────────────────────
\instruction{PUSHACC}{pushacc}

\paragraph{PDF Spec.}
\texttt{PUSHACC $n$}: Push \texttt{accu}, then set \texttt{accu = sp[$n$]}
(where $n$ counts from the \emph{new} stack top).  Specialized forms
\texttt{PUSHACC1}--\texttt{PUSHACC7}.

\paragraph{OCaml \texttt{interp.c} (line 356).}
\csnippet{PUSHACC}

\paragraph{Rocq \texttt{Interp.v} (line 119).}
\rocqsnippet{PUSHACC}

\paragraph{Analysis.}
In the C code, \texttt{*-{}-sp = accu} pushes first, then \texttt{accu = sp[$n$]}
reads from the new stack.  Our Rocq code constructs
\texttt{new\_stack := accu :: stack} then does \texttt{nth\_error new\_stack $n$},
which is the same indexing.  Exact match.

\paragraph{Status.} \statusok

\paragraph{Risk.} \risklow


% ─── POP ──────────────────────────────────────────────────────────────────────
\instruction{POP}{pop}

\paragraph{PDF Spec.}
\texttt{POP $n$}: Pops $n$ values from the stack (\texttt{sp += $n$}).

\paragraph{OCaml \texttt{interp.c} (line 363).}
\csnippet{POP}

\paragraph{Rocq \texttt{Interp.v} (line 126).}
\rocqsnippet{POP}

\paragraph{Analysis.}
\texttt{sp += $n$} corresponds to \texttt{skipn $n$ stack}.  Exact match.

\paragraph{Status.} \statusok

\paragraph{Risk.} \risklow


% ─── ASSIGN ───────────────────────────────────────────────────────────────────
\instruction{ASSIGN}{assign}

\paragraph{PDF Spec.}
\texttt{ASSIGN $n$}: Sets \texttt{sp[$n$] = accu}, then \texttt{accu = ()}.

\paragraph{OCaml \texttt{interp.c} (line 366).}
\csnippet{ASSIGN}

\paragraph{Rocq \texttt{Interp.v} (line 129).}
\rocqsnippet{ASSIGN}

\paragraph{Analysis.}
The C code does \texttt{sp[*pc++] = accu; accu = Val\_unit}.  Our code uses
\texttt{set\_nth stack $n$ accu}, which is a functional update returning
\texttt{Some new\_stack} on success.  The error case covers out-of-bounds
which is undefined behavior in C.  Exact match on valid inputs.

\paragraph{Status.} \statusok

\paragraph{Risk.} \risklow


%═══════════════════════════════════════════════════════════════════════════════
\section{Environment Access}
\label{sec:env}
%═══════════════════════════════════════════════════════════════════════════════

% ─── ENVACC ───────────────────────────────────────────────────────────────────
\instruction{ENVACC}{envacc}

\paragraph{PDF Spec.}
\texttt{ENVACC $n$}: Sets \texttt{accu = Field(env, $n$)}.  Specialized
variants \texttt{ENVACC1}--\texttt{ENVACC4}.

\paragraph{OCaml \texttt{interp.c} (line 373).}
\csnippet{ENVACC}

\paragraph{Rocq \texttt{Interp.v} (line 135).}
\rocqsnippet{ENVACC}

\paragraph{Analysis.}
\texttt{Field(env, $n$)} corresponds to \texttt{field\_or\_heap s env $n$}.
The \texttt{field\_or\_heap} helper handles both inline \texttt{Val\_block}
and heap-allocated \texttt{Val\_ptr}/\texttt{Val\_closure} uniformly,
accounting for the offset in closures.  Exact match.

\paragraph{Status.} \statusok

\paragraph{Risk.} \risklow


% ─── PUSHENVACC ───────────────────────────────────────────────────────────────
\instruction{PUSHENVACC}{pushenvacc}

\paragraph{PDF Spec.}
\texttt{PUSHENVACC $n$}: Push accu, then set \texttt{accu = Field(env, $n$)}.

\paragraph{OCaml \texttt{interp.c} (line 391).}
\csnippet{PUSHENVACC}

\paragraph{Rocq \texttt{Interp.v} (line 141).}
\rocqsnippet{PUSHENVACC}

\paragraph{Analysis.}
Same pattern as \texttt{PUSH} + \texttt{ENVACC}.  Exact match.

\paragraph{Status.} \statusok

\paragraph{Risk.} \risklow


% ─── OFFSETCLOSURE ────────────────────────────────────────────────────────────
\instruction{OFFSETCLOSURE}{offsetclosure}

\paragraph{PDF Spec.}
\texttt{OFFSETCLOSURE $n$}: Sets \texttt{accu = env + $n$ * sizeof(value)}.
This navigates within a flat closure block to access a different function
in a mutually recursive group.  Specialized: \texttt{OFFSETCLOSURE0} ($n=0$),
\texttt{OFFSETCLOSURE3} ($n=3$), \texttt{OFFSETCLOSUREM3} ($n=-3$).

\paragraph{OCaml \texttt{interp.c} (line 615).}
\csnippet{OFFSETCLOSURE}

\paragraph{Rocq \texttt{Interp.v} (line 362).}
\rocqsnippet{OFFSETCLOSURE}

\paragraph{Analysis.}
The C code uses pointer arithmetic: \texttt{accu = env + *pc * sizeof(value)}.
Our code adjusts the offset within a \texttt{Val\_closure}: the new offset
is \texttt{base\_ofs + ofs}.  The multiplication by \texttt{sizeof(value)} is
implicit because our offsets count fields, not bytes.

For \texttt{Val\_block} with offset 0, we return \texttt{env} itself, which is
correct for the self-reference case (\texttt{OFFSETCLOSURE0}).  Non-zero offsets
on non-closure values produce an error, which is safe since well-formed bytecode
always has a closure in \texttt{env} when \texttt{OFFSETCLOSURE} is executed.

\paragraph{Status.} \statusok

\paragraph{Risk.} \risklow


% ─── PUSHOFFSETCLOSURE ────────────────────────────────────────────────────────
\instruction{PUSHOFFSETCLOSURE}{pushoffsetclosure}

\paragraph{PDF Spec.}
\texttt{PUSHOFFSETCLOSURE $n$}: Push accu, then \texttt{OFFSETCLOSURE $n$}.

\paragraph{OCaml \texttt{interp.c} (line 613).}
\csnippet{PUSHOFFSETCLOSURE}

\paragraph{Rocq \texttt{Interp.v} (line 374).}
\rocqsnippet{PUSHOFFSETCLOSURE}

\paragraph{Analysis.}
Same as \texttt{PUSH} + \texttt{OFFSETCLOSURE}.  Exact match.

\paragraph{Status.} \statusok

\paragraph{Risk.} \risklow


%═══════════════════════════════════════════════════════════════════════════════
\section{Function Application}
\label{sec:apply}
%═══════════════════════════════════════════════════════════════════════════════

% ─── PUSH_RETADDR ─────────────────────────────────────────────────────────────
\instruction{PUSH\_RETADDR}{push_retaddr}

\paragraph{PDF Spec.}
Pushes a return frame \texttt{[pc+ofs, env, extra\_args]} onto the stack,
preparing for a multi-argument function call.

\paragraph{OCaml \texttt{interp.c} (line 400).}
\csnippet{PUSH_RETADDR}

\paragraph{Rocq \texttt{Interp.v} (line 150).}
\rocqsnippet{PUSH_RETADDR}

\paragraph{Analysis.}
The C code pushes \texttt{sp[0]=(value)(pc+*pc), sp[1]=env,
sp[2]=Val\_long(extra\_args)}.  Our code pushes
\texttt{Val\_int ret\_addr :: env :: Val\_int(Z.of\_nat extra\_args) :: stack}.
The return address is pre-decoded to an absolute offset in our model (the loader
resolves relative offsets), matching the C \texttt{pc + *pc}.  Exact match.

\paragraph{Status.} \statusok

\paragraph{Risk.} \risklow


% ─── APPLY ────────────────────────────────────────────────────────────────────
\instruction{APPLY}{apply}

\paragraph{PDF Spec.}
\texttt{APPLY $n$}: Direct call.  Sets \texttt{extra\_args = $n$ - 1},
\texttt{pc = Code\_val(accu)}, \texttt{env = accu}.  The $n$ arguments
are already on the stack (pushed by \texttt{PUSH\_RETADDR} sequence).

\paragraph{OCaml \texttt{interp.c} (line 408).}
\csnippet{APPLY}

\paragraph{Rocq \texttt{Interp.v} (line 155).}
\rocqsnippet{APPLY}

\paragraph{Analysis.}
The C code does \texttt{extra\_args = *pc - 1; pc = Code\_val(accu); env = accu}.
Our code extracts the code pointer via \texttt{get\_code\_ptr\_s} and sets
\texttt{env = accu} (the closure itself becomes the environment).
\texttt{extra\_args} is set to \texttt{n - 1}.  Exact match.

The C code also does \texttt{goto check\_stacks} for potential stack overflow
reallocation, which we skip (no stack overflow in our model).

\paragraph{Status.} \statusok

\paragraph{Risk.} \risklow


% ─── APPLY1 ───────────────────────────────────────────────────────────────────
\instruction{APPLY1}{apply1}

\paragraph{PDF Spec.}
\texttt{APPLY1}: Call with 1 argument.  Saves return frame inline:
\texttt{[arg1, pc, env, extra\_args]}, then jumps to closure.

\paragraph{OCaml \texttt{interp.c} (line 414).}
\csnippet{APPLY1}

\paragraph{Rocq \texttt{Interp.v} (line 163).}
\rocqsnippet{APPLY1}

\paragraph{Analysis.}
C: \texttt{sp[0]=arg1, sp[1]=pc, sp[2]=env, sp[3]=Val\_long(extra\_args)}.
Rocq: \texttt{arg1 :: Val\_int pc' :: env :: Val\_int(extra\_args) :: rest}.
Exact match.  \texttt{extra\_args = 0} in both cases.

\paragraph{Status.} \statusok

\paragraph{Risk.} \risklow


% ─── APPLY2 ───────────────────────────────────────────────────────────────────
\instruction{APPLY2}{apply2}

\paragraph{PDF Spec.}
\texttt{APPLY2}: Call with 2 arguments.

\paragraph{OCaml \texttt{interp.c} (line 426).}
\csnippet{APPLY2}

\paragraph{Rocq \texttt{Interp.v} (line 175).}
\rocqsnippet{APPLY2}

\paragraph{Analysis.}
Same pattern as \texttt{APPLY1} with 2 args.  \texttt{extra\_args = 1}.
Exact match.

\paragraph{Status.} \statusok

\paragraph{Risk.} \risklow


% ─── APPLY3 ───────────────────────────────────────────────────────────────────
\instruction{APPLY3}{apply3}

\paragraph{PDF Spec.}
\texttt{APPLY3}: Call with 3 arguments.

\paragraph{OCaml \texttt{interp.c} (line 440).}
\csnippet{APPLY3}

\paragraph{Rocq \texttt{Interp.v} (line 187).}
\rocqsnippet{APPLY3}

\paragraph{Analysis.}
Same pattern with 3 args.  \texttt{extra\_args = 2}.  Exact match.

\paragraph{Status.} \statusok

\paragraph{Risk.} \risklow


% ─── APPTERM ──────────────────────────────────────────────────────────────────
\instruction{APPTERM}{appterm}

\paragraph{PDF Spec.}
\texttt{APPTERM $n$ $s$}: Tail call.  Slides the top $n$ stack words down
over a frame of size $s$, then jumps to the closure in \texttt{accu}.
\texttt{extra\_args += $n$ - 1}.

\paragraph{OCaml \texttt{interp.c} (line 457).}
\csnippet{APPTERM}

\paragraph{Rocq \texttt{Interp.v} (line 200).}
\rocqsnippet{APPTERM}

\paragraph{Analysis.}
C: slides \texttt{sp[0..n-1]} to \texttt{sp+slotsize-n} via loop.
Rocq: \texttt{firstn nargs stack ++ skipn slotsize stack}.
Both produce the same result: the top $n$ values above the discarded
$slotsize - n$ values.  Exact match.

\paragraph{Status.} \statusok

\paragraph{Risk.} \risklow


% ─── APPTERM1 ─────────────────────────────────────────────────────────────────
\instruction{APPTERM1}{appterm1}

\paragraph{PDF Spec.}
\texttt{APPTERM1 $s$}: Tail call with 1 argument.

\paragraph{OCaml \texttt{interp.c} (line 472).}
\csnippet{APPTERM1}

\paragraph{Rocq \texttt{Interp.v} (line 209).}
\rocqsnippet{APPTERM1}

\paragraph{Analysis.}
C: \texttt{sp = sp + *pc - 1; sp[0] = arg1}.
Rocq: \texttt{arg1 :: skipn slotsize stack}.
Both keep 1 arg and discard the frame.  \texttt{extra\_args} unchanged.  Exact match.

\paragraph{Status.} \statusok

\paragraph{Risk.} \risklow


% ─── APPTERM2 ─────────────────────────────────────────────────────────────────
\instruction{APPTERM2}{appterm2}

\paragraph{PDF Spec.}
\texttt{APPTERM2 $s$}: Tail call with 2 arguments.

\paragraph{OCaml \texttt{interp.c} (line 480).}
\csnippet{APPTERM2}

\paragraph{Rocq \texttt{Interp.v} (line 221).}
\rocqsnippet{APPTERM2}

\paragraph{Analysis.}
Same pattern, 2 args.  \texttt{extra\_args += 1}.  Exact match.

\paragraph{Status.} \statusok

\paragraph{Risk.} \risklow


% ─── APPTERM3 ─────────────────────────────────────────────────────────────────
\instruction{APPTERM3}{appterm3}

\paragraph{PDF Spec.}
\texttt{APPTERM3 $s$}: Tail call with 3 arguments.

\paragraph{OCaml \texttt{interp.c} (line 491).}
\csnippet{APPTERM3}

\paragraph{Rocq \texttt{Interp.v} (line 233).}
\rocqsnippet{APPTERM3}

\paragraph{Analysis.}
Same pattern, 3 args.  \texttt{extra\_args += 2}.  Exact match.

\paragraph{Status.} \statusok

\paragraph{Risk.} \risklow


% ─── RETURN ───────────────────────────────────────────────────────────────────
\instruction{RETURN}{return}

\paragraph{PDF Spec.}
\texttt{RETURN $n$}: Pop $n$ locals.  If \texttt{extra\_args > 0}, decrement
and tail-call \texttt{accu}; otherwise restore the return frame
\texttt{[pc, env, extra\_args]}.

\paragraph{OCaml \texttt{interp.c} (line 505).}
\csnippet{RETURN}

\paragraph{Rocq \texttt{Interp.v} (line 246).}
\rocqsnippet{RETURN}

\paragraph{Analysis.}
C: \texttt{sp += *pc; if (extra\_args > 0) \{ extra\_args-{}-; pc = Code\_val(accu);
env = accu \} else \{ pc = sp[0]; env = sp[1]; extra\_args = Long\_val(sp[2]);
sp += 3 \}}.

Rocq: \texttt{skipn stacksize stack}, then checks \texttt{Nat.ltb 0 extra\_args}.
If true, decrements and calls closure.  Otherwise, pattern-matches the return
frame \texttt{Val\_int ret\_pc :: saved\_env :: Val\_int saved\_ea :: rest}.

The C check is \texttt{extra\_args > 0} while Rocq uses \texttt{Nat.ltb 0 extra\_args},
which is equivalent.  Exact match.

\paragraph{Status.} \statusok

\paragraph{Risk.} \risklow


%═══════════════════════════════════════════════════════════════════════════════
\section{Closure Creation}
\label{sec:closure}
%═══════════════════════════════════════════════════════════════════════════════

% ─── RESTART ──────────────────────────────────────────────────────────────────
\instruction{RESTART}{restart}

\paragraph{PDF Spec.}
\texttt{RESTART}: Restores partially applied arguments from the current closure
environment.  The closure layout is \texttt{[code, closinfo, saved\_env, arg0,
arg1, ...]}.  Extra args from field 3 onward are pushed back onto the stack,
and \texttt{env} is restored from field 2.

\paragraph{OCaml \texttt{interp.c} (line 520).}
\csnippet{RESTART}

\paragraph{Rocq \texttt{Interp.v} (line 262).}
\rocqsnippet{RESTART}

\paragraph{Analysis.}
C: \texttt{num\_args = Wosize\_val(env) - 3}; pushes \texttt{Field(env, i+3)}
for $i=0..n-1$; sets \texttt{env = Field(env, 2)}.

Rocq: for \texttt{Val\_closure addr ofs}, reads the heap block, computes
\texttt{fields := skipn ofs all\_fields}, then \texttt{num\_args = length(fields) - 3},
pushes \texttt{skipn 3 fields}, and sets env to \texttt{fields[2]}.

Key difference: the Rocq version handles the \texttt{Val\_closure} offset
correctly, accounting for infix closures within a shared block.  The C version
uses \texttt{Wosize\_val(env)} which gives the total block size, and
\texttt{env} already points to the right sub-closure.  In the Rocq model, the
offset adjusts the view.  Functionally equivalent.

\paragraph{Status.} \statusok

\paragraph{Risk.} \riskmed --- The closure offset arithmetic for mutually recursive functions needs careful testing.


% ─── GRAB ─────────────────────────────────────────────────────────────────────
\instruction{GRAB}{grab}

\paragraph{PDF Spec.}
\texttt{GRAB $n$}: If \texttt{extra\_args >= $n$}, consume them and continue.
Otherwise, build a partial application closure containing the saved arguments
and return to the caller.

\paragraph{OCaml \texttt{interp.c} (line 530).}
\csnippet{GRAB}

\paragraph{Rocq \texttt{Interp.v} (line 289).}
\rocqsnippet{GRAB}

\paragraph{Analysis.}
\textbf{Sufficient args path:} Both decrement \texttt{extra\_args} by
\texttt{required} and continue.  Exact match.

\textbf{Partial application path:}
\begin{itemize}
  \item C: Allocates closure of size \texttt{num\_args + 3} with tag
    \texttt{Closure\_tag}.  Layout: \texttt{Code\_val = pc - 3} (points to
    \texttt{RESTART}), \texttt{Closinfo\_val = Make\_closinfo(0, 2)},
    \texttt{Field(2) = env}, \texttt{Field(3..) = args}.
  \item Rocq: Allocates with fields \texttt{Val\_int(pc-1) :: closinfo :: env
    :: saved\_args}.  \texttt{closinfo = Val\_int 0}.
\end{itemize}

\textbf{Discrepancy: closinfo value.}  The C code uses
\texttt{Make\_closinfo(0, 2)} which encodes the arity and start-of-environment
offset as a tagged value.  Our code uses \texttt{Val\_int 0}.  This is
\textbf{safe} as long as nothing in our model inspects the closinfo field ---
and currently nothing does.  However, if we later need to introspect closure
metadata (e.g., for marshaling or debugging), this will need to be fixed.

\textbf{Discrepancy: code pointer offset.}  The C code points to \texttt{pc - 3}
(the \texttt{RESTART} instruction, which is 3 words before \texttt{GRAB}).
Our code uses \texttt{pc - 1}.  This difference exists because the C
instruction stream has the \texttt{RESTART} opcode at \texttt{pc-3} (accounting
for the GRAB opcode word and its operand word and then the RESTART before that),
while in our model, the \texttt{RESTART} instruction is at index
\texttt{pc - 1} in the instruction array (the instruction before \texttt{GRAB}).
Both correctly point to \texttt{RESTART}.

\paragraph{Status.} \statusneedstest{Verify closinfo=0 does not cause issues with RESTART; verify pc-1 correctly targets RESTART in all test programs.}

\paragraph{Risk.} \riskhigh{The closinfo field is hardcoded to 0 instead of Make\_closinfo(0, 2).  Safe currently but fragile.}


% ─── CLOSURE ──────────────────────────────────────────────────────────────────
\instruction{CLOSURE}{closure}

\paragraph{PDF Spec.}
\texttt{CLOSURE $n$ $ofs$}: Build a closure with $n$ environment variables
and code pointer at \texttt{pc + $ofs$}.  Layout:
\texttt{[code\_ptr, closinfo, v0, ..., v\_{n-1}]}.
If $n > 0$, push \texttt{accu} first (it becomes \texttt{v0}).

\paragraph{OCaml \texttt{interp.c} (line 551).}
\csnippet{CLOSURE}

\paragraph{Rocq \texttt{Interp.v} (line 310).}
\rocqsnippet{CLOSURE}

\paragraph{Analysis.}
Both conditionally push \texttt{accu} when \texttt{nvars > 0}, allocate a
block of \texttt{2 + nvars} fields with \texttt{Closure\_tag}, fill code pointer
and environment variables, and set \texttt{closinfo}.

Same closinfo discrepancy as \texttt{GRAB}: we use \texttt{Val\_int 0} vs.\
\texttt{Make\_closinfo(0, 2)}.

The code offset is pre-resolved in our model (absolute), matching the C
\texttt{pc + *pc} after loader processing.

\paragraph{Status.} \statusok

\paragraph{Risk.} \riskmed --- closinfo=0 (same systemic issue as GRAB).


% ─── CLOSUREREC ───────────────────────────────────────────────────────────────
\instruction{CLOSUREREC}{closurerec}

\paragraph{PDF Spec.}
\texttt{CLOSUREREC $nf$ $nv$ [$ofs_0$, ..., $ofs_{nf-1}$]}: Build a flat
mutual recursion closure block.  Total size is $3 \times nf - 1 + nv$.
Layout: \texttt{[code0, ci0, infix1, code1, ci1, ..., v0, ..., v\_{nv-1}]}.
Each function $i$ is accessed at offset $3i$ from the block start.
The closures are pushed onto the stack in order, with \texttt{accu} set to
closure 0.

\paragraph{OCaml \texttt{interp.c} (line 575).}
\csnippet{CLOSUREREC}

\paragraph{Rocq \texttt{Interp.v} (line 324).}
\rocqsnippet{CLOSUREREC}

\paragraph{Analysis.}
This is one of the most complex instructions and deserves careful attention.

\textbf{Block layout:}
\begin{itemize}
  \item C: Block size = $nf \times 3 - 1 + nv$.  Fields: \texttt{code0, closinfo0,
    [infix\_hdr, code\_i, closinfo\_i]\_{i=1..nf-1}, v0, ..., v\_{nv-1}}.
  \item Rocq: \texttt{build\_closure\_fields} produces the same layout:
    first function gets \texttt{code :: closinfo}, subsequent get
    \texttt{infix\_hdr :: code :: closinfo}, then environment variables.
\end{itemize}

\textbf{Infix headers:}
\begin{itemize}
  \item C: \texttt{Make\_header(i*3, Infix\_tag, Caml\_white)} --- this is a
    real GC header word encoding the infix offset.
  \item Rocq: \texttt{Val\_block Infix\_tag []} --- an empty block tagged
    \texttt{Infix\_tag}.  This does \textbf{not} encode the offset.
\end{itemize}

This is a \textbf{potential issue} if anything inspects the infix header to
compute offsets.  In practice, the header is only used by the GC (to scan
through the block), and our model has no GC.  The \texttt{Val\_closure addr ofs}
representation tracks offsets explicitly, so the infix header content is irrelevant
for execution.

\textbf{Closinfo:}  Same systemic issue: \texttt{Val\_int 0} vs.\
\texttt{Make\_closinfo(0, envofs)}.  The C version stores decreasing envofs
values for each sub-closure; we store 0.

\textbf{Push order:}  C pushes \texttt{accu} first (before allocation),
then pushes sub-closures \texttt{1..nf-1} during field initialization.
Rocq pushes via \texttt{push\_closures} after allocation.  Both result in
closure $nf-1$ on top and closure 0 in \texttt{accu}.  The push
function is recursive with \texttt{closure\_at i'} for $i' = 0..nf-1$,
pushing them bottom-up.

\paragraph{Status.} \statusneedstest{Verify push order matches for nf > 2; verify infix header irrelevance.}

\paragraph{Risk.} \riskhigh{Complex instruction with three systemic discrepancies (infix headers, closinfo, push ordering).  Covered by PBT but needs scrutiny.}


%═══════════════════════════════════════════════════════════════════════════════
\section{Global Variables}
\label{sec:globals}
%═══════════════════════════════════════════════════════════════════════════════

% ─── GETGLOBAL ────────────────────────────────────────────────────────────────
\instruction{GETGLOBAL}{getglobal}

\paragraph{PDF Spec.}
\texttt{GETGLOBAL $n$}: Sets \texttt{accu = Field(global\_data, $n$)}.

\paragraph{OCaml \texttt{interp.c} (line 637).}
\csnippet{GETGLOBAL}

\paragraph{Rocq \texttt{Interp.v} (line 387).}
\rocqsnippet{GETGLOBAL}

\paragraph{Analysis.}
C: \texttt{Field(caml\_global\_data, *pc)}.
Rocq: \texttt{nth\_error global $n$}.  Direct correspondence.

\paragraph{Status.} \statusok

\paragraph{Risk.} \risklow


% ─── PUSHGETGLOBAL ────────────────────────────────────────────────────────────
\instruction{PUSHGETGLOBAL}{pushgetglobal}

\paragraph{PDF Spec.}
\texttt{PUSHGETGLOBAL $n$}: Push accu, then \texttt{GETGLOBAL $n$}.

\paragraph{OCaml \texttt{interp.c} (line 634).}
\csnippet{PUSHGETGLOBAL}

\paragraph{Rocq \texttt{Interp.v} (line 393).}
\rocqsnippet{PUSHGETGLOBAL}

\paragraph{Analysis.}
Push + getglobal.  Exact match.

\paragraph{Status.} \statusok

\paragraph{Risk.} \risklow


% ─── GETGLOBALFIELD ───────────────────────────────────────────────────────────
\instruction{GETGLOBALFIELD}{getglobalfield}

\paragraph{PDF Spec.}
\texttt{GETGLOBALFIELD $n$ $p$}: \texttt{accu = Field(Field(global\_data, $n$), $p$)}.

\paragraph{OCaml \texttt{interp.c} (line 645).}
\csnippet{GETGLOBALFIELD}

\paragraph{Rocq \texttt{Interp.v} (line 400).}
\rocqsnippet{GETGLOBALFIELD}

\paragraph{Analysis.}
Two-level field access.  Rocq uses \texttt{nth\_error global $n$} then
\texttt{field\_or\_heap} for field $p$.  Exact match.

\paragraph{Status.} \statusok

\paragraph{Risk.} \risklow


% ─── PUSHGETGLOBALFIELD ───────────────────────────────────────────────────────
\instruction{PUSHGETGLOBALFIELD}{pushgetglobalfield}

\paragraph{PDF Spec.}
Push accu, then \texttt{GETGLOBALFIELD}.

\paragraph{OCaml \texttt{interp.c} (line 642).}
\csnippet{PUSHGETGLOBALFIELD}

\paragraph{Rocq \texttt{Interp.v} (line 410).}
\rocqsnippet{PUSHGETGLOBALFIELD}

\paragraph{Analysis.}
Same as push + getglobalfield.  Exact match.

\paragraph{Status.} \statusok

\paragraph{Risk.} \risklow


% ─── SETGLOBAL ────────────────────────────────────────────────────────────────
\instruction{SETGLOBAL}{setglobal}

\paragraph{PDF Spec.}
\texttt{SETGLOBAL $n$}: \texttt{Field(global\_data, $n$) = accu; accu = ()}.

\paragraph{OCaml \texttt{interp.c} (line 653).}
\csnippet{SETGLOBAL}

\paragraph{Rocq \texttt{Interp.v} (line 421).}
\rocqsnippet{SETGLOBAL}

\paragraph{Analysis.}
C: \texttt{caml\_modify(\&Field(caml\_global\_data, *pc), accu); accu = Val\_unit}.
Rocq: \texttt{set\_nth global $n$ accu}, falls back to unchanged on error (silently).
The silent fallback is slightly different from C (which would segfault on
out-of-bounds), but both cases are unreachable with well-formed bytecode.

\paragraph{Status.} \statusok

\paragraph{Risk.} \risklow


%═══════════════════════════════════════════════════════════════════════════════
\section{Allocation}
\label{sec:alloc}
%═══════════════════════════════════════════════════════════════════════════════

% ─── ATOM ─────────────────────────────────────────────────────────────────────
\instruction{ATOM}{atom}

\paragraph{PDF Spec.}
\texttt{ATOM $t$}: Sets \texttt{accu} to the atom with tag $t$ (a zero-size block).

\paragraph{OCaml \texttt{interp.c} (line 670).}
\csnippet{ATOM}

\paragraph{Rocq \texttt{Interp.v} (line 426).}
\rocqsnippet{ATOM}

\paragraph{Analysis.}
C: \texttt{Atom(*pc++)} returns a pointer to a pre-allocated zero-size block.
Rocq: \texttt{heap\_alloc s $t$ []} allocates a fresh empty block.  Semantically
equivalent---both produce a zero-size block with the given tag.  The Rocq version
allocates a new block each time (no sharing), which is safe since atoms are
immutable.

\paragraph{Status.} \statusok

\paragraph{Risk.} \risklow


% ─── PUSHATOM ─────────────────────────────────────────────────────────────────
\instruction{PUSHATOM}{pushatom}

\paragraph{PDF Spec.}
Push accu, then \texttt{ATOM}.

\paragraph{OCaml \texttt{interp.c} (line 667).}
\csnippet{PUSHATOM}

\paragraph{Rocq \texttt{Interp.v} (line 430).}
\rocqsnippet{PUSHATOM}

\paragraph{Analysis.}
Push + atom.  Exact match.

\paragraph{Status.} \statusok

\paragraph{Risk.} \risklow


% ─── MAKEBLOCK ────────────────────────────────────────────────────────────────
\instruction{MAKEBLOCK}{makeblock}

\paragraph{PDF Spec.}
\texttt{MAKEBLOCK $s$ $t$}: Allocate a block of size $s$ with tag $t$.
\texttt{Field(0) = accu}, remaining fields popped from stack.

\paragraph{OCaml \texttt{interp.c} (line 673).}
\csnippet{MAKEBLOCK}

\paragraph{Rocq \texttt{Interp.v} (line 436).}
\rocqsnippet{MAKEBLOCK}

\paragraph{Analysis.}
C: \texttt{Field(block, 0) = accu; for (i=1; i<wosize; i++) Field(block,i) = *sp++}.
Rocq: \texttt{accu :: firstn (size-1) stack}, then heap\_alloc.

Note the C operand order is \texttt{wosize, tag} while our AST stores
\texttt{tag, size} --- the loader handles this reordering.

Exact match on semantics.

\paragraph{Status.} \statusok

\paragraph{Risk.} \risklow


% ─── MAKEBLOCK1 ───────────────────────────────────────────────────────────────
\instruction{MAKEBLOCK1}{makeblock1}

\paragraph{PDF Spec.}
\texttt{MAKEBLOCK1 $t$}: 1-field block.

\paragraph{OCaml \texttt{interp.c} (line 690).}
\csnippet{MAKEBLOCK1}

\paragraph{Rocq \texttt{Interp.v} (line 442).}
\rocqsnippet{MAKEBLOCK1}

\paragraph{Analysis.}
Single-field block with accu.  Exact match.

\paragraph{Status.} \statusok

\paragraph{Risk.} \risklow


% ─── MAKEBLOCK2 ───────────────────────────────────────────────────────────────
\instruction{MAKEBLOCK2}{makeblock2}

\paragraph{PDF Spec.}
\texttt{MAKEBLOCK2 $t$}: 2-field block.

\paragraph{OCaml \texttt{interp.c} (line 698).}
\csnippet{MAKEBLOCK2}

\paragraph{Rocq \texttt{Interp.v} (line 446).}
\rocqsnippet{MAKEBLOCK2}

\paragraph{Analysis.}
\texttt{[accu, sp[0]]}.  Exact match.

\paragraph{Status.} \statusok

\paragraph{Risk.} \risklow


% ─── MAKEBLOCK3 ───────────────────────────────────────────────────────────────
\instruction{MAKEBLOCK3}{makeblock3}

\paragraph{PDF Spec.}
\texttt{MAKEBLOCK3 $t$}: 3-field block.

\paragraph{OCaml \texttt{interp.c} (line 708).}
\csnippet{MAKEBLOCK3}

\paragraph{Rocq \texttt{Interp.v} (line 454).}
\rocqsnippet{MAKEBLOCK3}

\paragraph{Analysis.}
\texttt{[accu, sp[0], sp[1]]}.  Exact match.

\paragraph{Status.} \statusok

\paragraph{Risk.} \risklow


% ─── MAKEFLOATBLOCK ───────────────────────────────────────────────────────────
\instruction{MAKEFLOATBLOCK}{makefloatblock}

\paragraph{PDF Spec.}
\texttt{MAKEFLOATBLOCK $n$}: Allocates a block of $n$ boxed floats with tag
\texttt{Double\_array\_tag} (254).

\paragraph{OCaml \texttt{interp.c} (line 719).}
\csnippet{MAKEFLOATBLOCK}

\paragraph{Rocq \texttt{Interp.v} (line 464).}
\rocqsnippet{MAKEFLOATBLOCK}

\paragraph{Analysis.}
C allocates with \texttt{size * Double\_wosize} words and uses
\texttt{Store\_double\_flat\_field} for flat (unboxed) float storage.
Rocq allocates with tag 254 and stores values directly (no float unboxing).
This is a simplification: our model does not distinguish boxed from unboxed
float representation.  Since OCaml bytecode programs that use floats will
have their values go through C primitives for arithmetic, and we model C calls
via \texttt{CCall\_request}, this simplification is acceptable for integer-only
programs but would need revision for float-heavy code.

\paragraph{Status.} \statusneedstest{Float representation differs from C; needs testing with float-using programs.}

\paragraph{Risk.} \riskmed --- Simplified float model.


%═══════════════════════════════════════════════════════════════════════════════
\section{Field Access}
\label{sec:field}
%═══════════════════════════════════════════════════════════════════════════════

% ─── GETFIELD ─────────────────────────────────────────────────────────────────
\instruction{GETFIELD}{getfield}

\paragraph{PDF Spec.}
\texttt{GETFIELD $n$}: \texttt{accu = Field(accu, $n$)}.

\paragraph{OCaml \texttt{interp.c} (line 747).}
\csnippet{GETFIELD}

\paragraph{Rocq \texttt{Interp.v} (line 470).}
\rocqsnippet{GETFIELD}

\paragraph{Analysis.}
\texttt{field\_or\_heap s accu $n$}.  Handles both inline blocks and heap pointers.
Exact match.

\paragraph{Status.} \statusok

\paragraph{Risk.} \risklow


% ─── GETFLOATFIELD ────────────────────────────────────────────────────────────
\instruction{GETFLOATFIELD}{getfloatfield}

\paragraph{PDF Spec.}
\texttt{GETFLOATFIELD $n$}: Reads the $n$-th double from a float array block
and boxes it.

\paragraph{OCaml \texttt{interp.c} (line 749).}
\csnippet{GETFLOATFIELD}

\paragraph{Rocq \texttt{Interp.v} (line 477).}
\rocqsnippet{GETFLOATFIELD}

\paragraph{Analysis.}
C: reads flat double, allocates new box (\texttt{Double\_tag}).
Rocq: just reads the field (no boxing/unboxing distinction).
Same simplification as \texttt{MAKEFLOATBLOCK}.

\paragraph{Status.} \statusneedstest{Float boxing semantics simplified.}

\paragraph{Risk.} \riskmed --- Float model simplification.


% ─── SETFIELD ─────────────────────────────────────────────────────────────────
\instruction{SETFIELD}{setfield}

\paragraph{PDF Spec.}
\texttt{SETFIELD $n$}: \texttt{Field(accu, $n$) = sp[0]; sp++; accu = ()}.
Mutates field $n$ of the block in \texttt{accu} with the value from stack top.

\paragraph{OCaml \texttt{interp.c} (line 772).}
\csnippet{SETFIELD}

\paragraph{Rocq \texttt{Interp.v} (line 483).}
\rocqsnippet{SETFIELD}

\paragraph{Analysis.}
C: \texttt{caml\_modify(\&Field(accu, *pc), *sp++); accu = Val\_unit}.
Rocq: pops new value from stack, looks up heap block, uses \texttt{set\_nth}
on fields, then \texttt{heap\_update}.

\textbf{Important:} In the C code, \texttt{accu} is the block and \texttt{*sp}
is the new value.  In the Rocq code (line 484), the new value comes from the
stack (\texttt{newval :: rest}), and \texttt{accu} (\texttt{Val\_ptr addr}) is
the block.  This matches the C semantics correctly.

\paragraph{Status.} \statusok

\paragraph{Risk.} \risklow


% ─── SETFLOATFIELD ────────────────────────────────────────────────────────────
\instruction{SETFLOATFIELD}{setfloatfield}

\paragraph{PDF Spec.}
\texttt{SETFLOATFIELD $n$}: Stores a float into field $n$ of a float array.
\texttt{Store\_double\_flat\_field(accu, $n$, Double\_val(*sp)); sp++; accu = ()}.

\paragraph{OCaml \texttt{interp.c} (line 777).}
\csnippet{SETFLOATFIELD}

\paragraph{Rocq \texttt{Interp.v} (line 504).}
\rocqsnippet{SETFLOATFIELD}

\paragraph{Analysis.}
\textbf{CRITICAL: Potential operand inversion.}

In \texttt{interp.c} (line 778):
\texttt{Store\_double\_flat\_field(\textbf{accu}, *pc, Double\_val(*sp))}.
Here \texttt{accu} is the float array and \texttt{*sp} is the new float value.

In our Rocq code (line 504--521), the code does:
\begin{quote}
\texttt{match s.(stack) with | arr\_val :: rest => ...}
\end{quote}
So \texttt{arr\_val} (from the stack) is treated as the float array, and
\texttt{s.(accu)} is the new value.  This is the \textbf{opposite} convention
from \texttt{interp.c}!

However, looking more carefully at the comment on line 503:
``\textit{SETFLOATFIELD n: stack top is the float array (tag 254), accu is the new float value.}''
This comment explicitly states the inverted convention.

\textbf{Verdict:} The Rocq code has \texttt{accu} = new value and
\texttt{sp[0]} = array.  The C code has \texttt{accu} = array and
\texttt{sp[0]} = new value.  This is a \textbf{bug} if both conventions
are used with the same bytecode.

\textbf{BUT:} Looking at the regular \texttt{SETFIELD} (line 483), it correctly
uses accu = block and sp = new value.  \texttt{SETFLOATFIELD} should follow
the same convention.  The inversion in our Rocq code is indeed a bug.

\paragraph{Status.} \statusbug{Operand convention inverted: accu should be the float array, sp[0] should be the new value, but our code swaps them.}

\paragraph{Risk.} \riskhigh{Affects all programs using mutable float arrays.  May be masked by PBT if float programs are rare in the test suite.}


%═══════════════════════════════════════════════════════════════════════════════
\section{Array and String Operations}
\label{sec:array}
%═══════════════════════════════════════════════════════════════════════════════

% ─── VECTLENGTH ───────────────────────────────────────────────────────────────
\instruction{VECTLENGTH}{vectlength}

\paragraph{PDF Spec.}
\texttt{VECTLENGTH}: \texttt{accu = Wosize\_val(accu)}.  For double arrays,
divides by \texttt{Double\_wosize}.

\paragraph{OCaml \texttt{interp.c} (line 786).}
\csnippet{VECTLENGTH}

\paragraph{Rocq \texttt{Interp.v} (line 524).}
\rocqsnippet{VECTLENGTH}

\paragraph{Analysis.}
C: checks tag for \texttt{Double\_array\_tag} and divides size.
Rocq: uses \texttt{size\_or\_heap} which returns the number of fields directly.
Since our float arrays store one value per field (no flat representation),
the count is already correct without division.

\paragraph{Status.} \statusok

\paragraph{Risk.} \risklow


% ─── GETVECTITEM ──────────────────────────────────────────────────────────────
\instruction{GETVECTITEM}{getvectitem}

\paragraph{PDF Spec.}
\texttt{GETVECTITEM}: \texttt{accu = Field(accu, Long\_val(sp[0])); sp++}.

\paragraph{OCaml \texttt{interp.c} (line 795).}
\csnippet{GETVECTITEM}

\paragraph{Rocq \texttt{Interp.v} (line 530).}
\rocqsnippet{GETVECTITEM}

\paragraph{Analysis.}
Index from stack top, field access on accu.  Exact match.

\paragraph{Status.} \statusok

\paragraph{Risk.} \risklow


% ─── SETVECTITEM ──────────────────────────────────────────────────────────────
\instruction{SETVECTITEM}{setvectitem}

\paragraph{PDF Spec.}
\texttt{SETVECTITEM}: \texttt{Field(accu, Long\_val(sp[0])) = sp[1]; sp+=2; accu = ()}.

\paragraph{OCaml \texttt{interp.c} (line 799).}
\csnippet{SETVECTITEM}

\paragraph{Rocq \texttt{Interp.v} (line 540).}
\rocqsnippet{SETVECTITEM}

\paragraph{Analysis.}
C: \texttt{sp[0]} = index, \texttt{sp[1]} = new value.
Rocq: \texttt{Val\_int idx :: newval :: rest}.  Same order.
Mutates via \texttt{heap\_update}.  Exact match.

\paragraph{Status.} \statusok

\paragraph{Risk.} \risklow


% ─── GETSTRINGCHAR / GETBYTESCHAR ────────────────────────────────────────────
\instruction{GETSTRINGCHAR}{getstringchar}

\paragraph{PDF Spec.}
\texttt{GETSTRINGCHAR}/\texttt{GETBYTESCHAR}:
\texttt{accu = Val\_int(Byte\_u(accu, Long\_val(sp[0]))); sp++}.

\paragraph{OCaml \texttt{interp.c} (line 806).}
\csnippet{GETSTRINGCHAR}

\paragraph{Rocq \texttt{Interp.v} (line 560).}
\rocqsnippet{GETBYTESCHAR}

\paragraph{Analysis.}
Our model represents strings as blocks of integer character values (not raw bytes).
The code extracts a \texttt{Val\_int c} at the given index.  This is a
simplification---real OCaml strings are packed byte arrays.  The simplification
is valid as long as string operations go through C primitives and our loader
correctly unpacks strings.

Both \texttt{GETSTRINGCHAR} and \texttt{GETBYTESCHAR} share the same case
in our code, matching the C implementation.

\paragraph{Status.} \statusok

\paragraph{Risk.} \risklow


% ─── GETBYTESCHAR ─────────────────────────────────────────────────────────────
\instruction{GETBYTESCHAR}{getbyteschar}

\paragraph{PDF Spec.}
Same as \texttt{GETSTRINGCHAR}.

\paragraph{Analysis.}
Combined with \texttt{GETSTRINGCHAR} in both C and Rocq.

\paragraph{Status.} \statusok

\paragraph{Risk.} \risklow


% ─── SETBYTESCHAR ─────────────────────────────────────────────────────────────
\instruction{SETBYTESCHAR}{setbyteschar}

\paragraph{PDF Spec.}
\texttt{SETBYTESCHAR}: \texttt{Byte\_u(accu, Long\_val(sp[0])) = Int\_val(sp[1]); sp+=2; accu = ()}.

\paragraph{OCaml \texttt{interp.c} (line 811).}
\csnippet{SETBYTESCHAR}

\paragraph{Rocq \texttt{Interp.v} (line 570).}
\rocqsnippet{SETBYTESCHAR}

\paragraph{Analysis.}
Heap mutation for byte arrays.  Our model uses \texttt{set\_nth} on the field
list.  Functionally equivalent given our string-as-block-of-ints representation.

\paragraph{Status.} \statusok

\paragraph{Risk.} \risklow


%═══════════════════════════════════════════════════════════════════════════════
\section{Control Flow}
\label{sec:control}
%═══════════════════════════════════════════════════════════════════════════════

% ─── BRANCH ───────────────────────────────────────────────────────────────────
\instruction{BRANCH}{branch}

\paragraph{PDF Spec.}
\texttt{BRANCH $ofs$}: \texttt{pc += $ofs$}.

\paragraph{OCaml \texttt{interp.c} (line 819).}
\csnippet{BRANCH}

\paragraph{Rocq \texttt{Interp.v} (line 590).}
\rocqsnippet{BRANCH}

\paragraph{Analysis.}
C: \texttt{pc += *pc} (relative offset).
Rocq: \texttt{pc = target} (absolute, pre-resolved by loader).
Equivalent after loader processing.

\paragraph{Status.} \statusok

\paragraph{Risk.} \risklow


% ─── BRANCHIF ─────────────────────────────────────────────────────────────────
\instruction{BRANCHIF}{branchif}

\paragraph{PDF Spec.}
\texttt{BRANCHIF $ofs$}: Branch if \texttt{accu != Val\_false} (i.e., $\neq 0$).

\paragraph{OCaml \texttt{interp.c} (line 822).}
\csnippet{BRANCHIF}

\paragraph{Rocq \texttt{Interp.v} (line 593).}
\rocqsnippet{BRANCHIF}

\paragraph{Analysis.}
C: \texttt{accu != Val\_false}.  \texttt{Val\_false = Val\_int(0)}.
Rocq: matches \texttt{Val\_int 0} for the no-branch case, wildcard for branch.
This is correct: any non-\texttt{Val\_int 0} value (including blocks) triggers
the branch, matching the C \texttt{Is\_block} behavior.

\paragraph{Status.} \statusok

\paragraph{Risk.} \risklow


% ─── BRANCHIFNOT ──────────────────────────────────────────────────────────────
\instruction{BRANCHIFNOT}{branchifnot}

\paragraph{PDF Spec.}
\texttt{BRANCHIFNOT $ofs$}: Branch if \texttt{accu == Val\_false}.

\paragraph{OCaml \texttt{interp.c} (line 825).}
\csnippet{BRANCHIFNOT}

\paragraph{Rocq \texttt{Interp.v} (line 599).}
\rocqsnippet{BRANCHIFNOT}

\paragraph{Analysis.}
Dual of \texttt{BRANCHIF}.  Exact match.

\paragraph{Status.} \statusok

\paragraph{Risk.} \risklow


% ─── SWITCH ───────────────────────────────────────────────────────────────────
\instruction{SWITCH}{switch}

\paragraph{PDF Spec.}
\texttt{SWITCH}: Dispatches on \texttt{accu}.  The operand word encodes
\texttt{nc} (number of constant cases, low 16 bits) and \texttt{nb} (number
of block cases, high 16 bits).  If \texttt{accu} is an integer, jump to
\texttt{pc[Long\_val(accu)]}.  If a block, jump to \texttt{pc[nc + Tag\_val(accu)]}.

\paragraph{OCaml \texttt{interp.c} (line 828).}
\csnippet{SWITCH}

\paragraph{Rocq \texttt{Interp.v} (line 605).}
\rocqsnippet{SWITCH}

\paragraph{Analysis.}
C: \texttt{sizes = *pc++; if (Is\_block(accu)) index = Tag\_val(accu),
pc += pc[(sizes \& 0xFFFF) + index]; else index = Long\_val(accu),
pc += pc[index]}.

Rocq: The instruction is pre-decoded with separate \texttt{const\_targets}
and \texttt{block\_targets} lists.  For integers, looks up
\texttt{nth\_error const\_targets (Z.to\_nat n)}.  For blocks/pointers,
uses \texttt{tag\_or\_heap} to get the tag, then
\texttt{nth\_error block\_targets t}.

The key difference is that the loader splits the single offset table into
two separate lists, while \texttt{interp.c} uses a combined table with
\texttt{nc} as the split point.  Functionally equivalent assuming the
loader correctly splits at \texttt{nc}.

\paragraph{Status.} \statusok

\paragraph{Risk.} \risklow --- Depends on correct loader splitting.


% ─── BOOLNOT ──────────────────────────────────────────────────────────────────
\instruction{BOOLNOT}{boolnot}

\paragraph{PDF Spec.}
\texttt{BOOLNOT}: \texttt{accu = Val\_not(accu)}.

\paragraph{OCaml \texttt{interp.c} (line 841).}
\csnippet{BOOLNOT}

\paragraph{Rocq \texttt{Interp.v} (line 628).}
\rocqsnippet{BOOLNOT}

\paragraph{Analysis.}
C: \texttt{Val\_not(accu)} flips the tag bit.
Rocq: pattern matches on \texttt{Val\_int 0} to produce \texttt{val\_true}/\texttt{val\_false}.
Exact match for boolean values.

\paragraph{Status.} \statusok

\paragraph{Risk.} \risklow


%═══════════════════════════════════════════════════════════════════════════════
\section{Exception Handling}
\label{sec:exception}
%═══════════════════════════════════════════════════════════════════════════════

% ─── PUSHTRAP ─────────────────────────────────────────────────────────────────
\instruction{PUSHTRAP}{pushtrap}

\paragraph{PDF Spec.}
\texttt{PUSHTRAP $ofs$}: Push trap frame \texttt{[handler\_pc, trap\_link,
env, extra\_args]} and set \texttt{trapsp = sp}.

\paragraph{OCaml \texttt{interp.c} (line 847).}
\csnippet{PUSHTRAP}

\paragraph{Rocq \texttt{Interp.v} (line 637).}
\rocqsnippet{PUSHTRAP}

\paragraph{Analysis.}
C trap frame: \texttt{sp[0] = pc+*pc} (handler), \texttt{sp[1] =
Val\_long(trapsp - sp)} (link as offset), \texttt{sp[2] = env},
\texttt{sp[3] = Val\_long(extra\_args)}.  Then \texttt{trapsp = sp}.

Rocq trap frame: \texttt{Val\_int handler\_pc :: Val\_int(trap\_sp) ::
env :: Val\_int(extra\_args) :: stack}.  Then \texttt{trap\_sp = length(new\_stack)}.

\textbf{Trap link encoding:}
\begin{itemize}
  \item C stores \texttt{trapsp - sp} --- the \emph{offset} from current sp
    to previous trap frame.
  \item Rocq stores the \emph{absolute} previous \texttt{trap\_sp} value
    (stack depth).
\end{itemize}

These are equivalent representations.  The C offset plus the current sp gives
the previous trapsp.  Our absolute value directly records it.  On
\texttt{POPTRAP}, the C code recovers \texttt{trapsp = sp +
Long\_val(Trap\_link\_offset(sp))}, while we directly restore the saved
\texttt{trap\_sp} value.

\paragraph{Status.} \statusok

\paragraph{Risk.} \riskmed --- The trap link representation differs; any bug in the offset computation would be subtle.


% ─── POPTRAP ──────────────────────────────────────────────────────────────────
\instruction{POPTRAP}{poptrap}

\paragraph{PDF Spec.}
\texttt{POPTRAP}: Restore \texttt{trapsp} from the trap link and pop the 4-word
trap frame.

\paragraph{OCaml \texttt{interp.c} (line 857).}
\csnippet{POPTRAP}

\paragraph{Rocq \texttt{Interp.v} (line 645).}
\rocqsnippet{POPTRAP}

\paragraph{Analysis.}
C: \texttt{trapsp = sp + Long\_val(Trap\_link\_offset(sp)); sp += 4}.
Also checks \texttt{caml\_something\_to\_do} (signal handling) which we skip.

Rocq: Pattern-matches the trap frame, extracts \texttt{prev\_tsp} from
position 1, pops 4 values.  Sets \texttt{trap\_sp = Z.to\_nat prev\_tsp}.

Exact match modulo signal handling (which is a no-op in our model).

\paragraph{Status.} \statusok

\paragraph{Risk.} \risklow


% ─── RAISE ────────────────────────────────────────────────────────────────────
\instruction{RAISE}{raise}

\paragraph{PDF Spec.}
\texttt{RAISE}: Raise the exception in \texttt{accu}.  Restores sp to trapsp,
extracts handler pc, restores previous trap link, env, and extra\_args.

\paragraph{OCaml \texttt{interp.c} (line 881).}
\csnippet{RAISE}

\paragraph{Rocq \texttt{Interp.v} (line 655).}
\rocqsnippet{RAISE}

\paragraph{Analysis.}
All three raise variants (RAISE, RERAISE, RAISE\_NOTRACE) call \texttt{do\_raise}
in our code.  The C versions differ only in backtrace handling (which we don't model).

\texttt{do\_raise}: If \texttt{trap\_sp = 0}, returns
\texttt{Error "unhandled exception"}.  Otherwise computes
\texttt{k = length(stack) - trap\_sp}, skips $k$ elements to reach the trap
frame, extracts \texttt{[handler\_pc, prev\_tsp, saved\_env, saved\_ea]},
and jumps.

The C code: \texttt{sp = trapsp; pc = Trap\_pc(sp);
trapsp = sp + Long\_val(Trap\_link\_offset(sp)); env = sp[2];
extra\_args = Long\_val(sp[3]); sp += 4}.

Equivalent: our \texttt{skipn k stack} positions at the trap frame top, then
pattern-matches the 4-word frame.

\textbf{Edge case:} C checks if trapsp is beyond the initial stack (unhandled
exception in callback) via pointer comparison.  Our check
\texttt{trap\_sp = 0} is the equivalent sentinel for ``no handler''.

\paragraph{Status.} \statusok

\paragraph{Risk.} \riskmed --- Trap frame navigation is the most subtle part of the interpreter.


% ─── RERAISE ──────────────────────────────────────────────────────────────────
\instruction{RERAISE}{reraise}

\paragraph{PDF Spec.}
Same as \texttt{RAISE} but with backtrace re-raise semantics.

\paragraph{Analysis.}
Same as \texttt{RAISE} in our model (no backtrace).

\paragraph{Status.} \statusok

\paragraph{Risk.} \risklow


% ─── RAISE\_NOTRACE ───────────────────────────────────────────────────────────
\instruction{RAISE\_NOTRACE}{raise_notrace}

\paragraph{PDF Spec.}
Same as \texttt{RAISE} but without backtrace recording.

\paragraph{Analysis.}
Same as \texttt{RAISE} in our model.

\paragraph{Status.} \statusok

\paragraph{Risk.} \risklow


%═══════════════════════════════════════════════════════════════════════════════
\section{Signal Handling and C Calls}
\label{sec:ccall}
%═══════════════════════════════════════════════════════════════════════════════

% ─── CHECK_SIGNALS ────────────────────────────────────────────────────────────
\instruction{CHECK\_SIGNALS}{check_signals}

\paragraph{PDF Spec.}
\texttt{CHECK\_SIGNALS}: Process pending signals.

\paragraph{OCaml \texttt{interp.c} (line 916).}
\csnippet{CHECK_SIGNALS}

\paragraph{Rocq \texttt{Interp.v} (line 658).}
\rocqsnippet{CHECK_SIGNALS}

\paragraph{Analysis.}
Modeled as no-op (just advance PC).  In the C interpreter, this checks
\texttt{caml\_something\_to\_do} and processes pending actions.  Safe to skip
since we don't model OS signals.

\paragraph{Status.} \statusok

\paragraph{Risk.} \risklow


% ─── C_CALL ───────────────────────────────────────────────────────────────────
\instruction{C\_CALL}{c_call}

\paragraph{PDF Spec.}
\texttt{C\_CALL $n$ $p$}: Call C primitive $p$ with $n$ arguments.
The C code has specialized forms \texttt{C\_CALL1} through \texttt{C\_CALL5}
and a generic \texttt{C\_CALLN}.

\paragraph{OCaml \texttt{interp.c} (line 928).}
\csnippet{C_CALL}

\paragraph{Rocq \texttt{Interp.v} (line 662).}
\rocqsnippet{C_CALL}

\paragraph{Analysis.}
C: sets up frame, calls primitive, restores frame.
Rocq: returns \texttt{CCall\_request prim\_idx args cont}, suspending
execution.  The caller must supply the primitive's return value to resume.

This is a fundamental architectural difference --- we don't execute C primitives
but instead produce an interaction-tree-style request.  The PBT harness
fulfills these by calling the actual C primitives.

\paragraph{Status.} \statusok

\paragraph{Risk.} \risklow --- The interaction tree model is well-understood.


%═══════════════════════════════════════════════════════════════════════════════
\section{Integer Constants}
\label{sec:const}
%═══════════════════════════════════════════════════════════════════════════════

% ─── CONSTINT ─────────────────────────────────────────────────────────────────
\instruction{CONSTINT}{constint}

\paragraph{PDF Spec.}
\texttt{CONSTINT $n$}: \texttt{accu = Val\_int($n$)}.  Specialized:
\texttt{CONST0}--\texttt{CONST3}.

\paragraph{OCaml \texttt{interp.c} (line 996).}
\csnippet{CONSTINT}

\paragraph{Rocq \texttt{Interp.v} (line 668).}
\rocqsnippet{CONSTINT}

\paragraph{Analysis.}
Direct assignment.  Exact match.

\paragraph{Status.} \statusok

\paragraph{Risk.} \risklow


% ─── PUSHCONSTINT ─────────────────────────────────────────────────────────────
\instruction{PUSHCONSTINT}{pushconstint}

\paragraph{PDF Spec.}
\texttt{PUSHCONSTINT $n$}: Push accu, then \texttt{CONSTINT}.

\paragraph{OCaml \texttt{interp.c} (line 993).}
\csnippet{PUSHCONSTINT}

\paragraph{Rocq \texttt{Interp.v} (line 671).}
\rocqsnippet{PUSHCONSTINT}

\paragraph{Analysis.}
Push + constint.  Exact match.

\paragraph{Status.} \statusok

\paragraph{Risk.} \risklow


%═══════════════════════════════════════════════════════════════════════════════
\section{Integer Arithmetic}
\label{sec:arith}
%═══════════════════════════════════════════════════════════════════════════════

% ─── NEGINT ───────────────────────────────────────────────────────────────────
\instruction{NEGINT}{negint}

\paragraph{PDF Spec.}
\texttt{NEGINT}: \texttt{accu = -accu} (integer negation).

\paragraph{OCaml \texttt{interp.c} (line 1003).}
\csnippet{NEGINT}

\paragraph{Rocq \texttt{Interp.v} (line 675).}
\rocqsnippet{NEGINT}

\paragraph{Analysis.}
C: \texttt{(value)(2 - (intnat)accu)} --- this works because for tagged integers,
\texttt{Val\_int(n) = 2n+1}, so \texttt{2 - (2n+1) = -(2n-1) = 2(-n)+1 = Val\_int(-n)}.
Rocq: \texttt{Val\_int(-n)}.  Exact match.

\paragraph{Status.} \statusok

\paragraph{Risk.} \risklow


% ─── ADDINT ───────────────────────────────────────────────────────────────────
\instruction{ADDINT}{addint}

\paragraph{PDF Spec.}
\texttt{ADDINT}: \texttt{accu = accu + sp[0]; sp++}.

\paragraph{OCaml \texttt{interp.c} (line 1005).}
\csnippet{ADDINT}

\paragraph{Rocq \texttt{Interp.v} (line 681).}
\rocqsnippet{ADDINT}

\paragraph{Analysis.}
C: \texttt{(intnat)accu + (intnat)*sp++ - 1} (tag bit arithmetic).
Rocq: \texttt{Val\_int(a + b)}.  Exact match at the semantic level.

\paragraph{Status.} \statusok

\paragraph{Risk.} \risklow


% ─── SUBINT ───────────────────────────────────────────────────────────────────
\instruction{SUBINT}{subint}

\paragraph{PDF Spec.}
\texttt{SUBINT}: \texttt{accu = accu - sp[0]; sp++}.

\paragraph{OCaml \texttt{interp.c} (line 1007).}
\csnippet{SUBINT}

\paragraph{Rocq \texttt{Interp.v} (line 687).}
\rocqsnippet{SUBINT}

\paragraph{Analysis.}
C: \texttt{(intnat)accu - (intnat)*sp++ + 1}.  Rocq: \texttt{a - b}.  Exact match.

\paragraph{Status.} \statusok

\paragraph{Risk.} \risklow


% ─── MULINT ───────────────────────────────────────────────────────────────────
\instruction{MULINT}{mulint}

\paragraph{PDF Spec.}
\texttt{MULINT}: \texttt{accu = accu * sp[0]; sp++}.

\paragraph{OCaml \texttt{interp.c} (line 1009).}
\csnippet{MULINT}

\paragraph{Rocq \texttt{Interp.v} (line 693).}
\rocqsnippet{MULINT}

\paragraph{Analysis.}
C: \texttt{Val\_long(Long\_val(accu) * Long\_val(*sp++))}.  Rocq: \texttt{a * b}.
Exact match.

\paragraph{Status.} \statusok

\paragraph{Risk.} \risklow


% ─── DIVINT ───────────────────────────────────────────────────────────────────
\instruction{DIVINT}{divint}

\paragraph{PDF Spec.}
\texttt{DIVINT}: \texttt{accu = accu / sp[0]; sp++}.  Raises
\texttt{Division\_by\_zero} if divisor is 0.

\paragraph{OCaml \texttt{interp.c} (line 1012).}
\csnippet{DIVINT}

\paragraph{Rocq \texttt{Interp.v} (line 699).}
\rocqsnippet{DIVINT}

\paragraph{Analysis.}
C: checks \texttt{divisor == 0}, calls \texttt{caml\_raise\_zero\_divide()}.
Rocq: checks \texttt{Z.eqb b 0}, calls \texttt{do\_raise div\_by\_zero\_exn}.

Both raise an exception on zero.  The Rocq exception is a manually constructed
\texttt{Division\_by\_zero} block (tag 248, id -6), matching OCaml's predefined
exception.

Division uses \texttt{Z.quot} (truncation toward zero), matching C's
\texttt{/} operator for signed integers.

\paragraph{Status.} \statusok

\paragraph{Risk.} \risklow


% ─── MODINT ───────────────────────────────────────────────────────────────────
\instruction{MODINT}{modint}

\paragraph{PDF Spec.}
\texttt{MODINT}: \texttt{accu = accu \% sp[0]; sp++}.  Raises on zero.

\paragraph{OCaml \texttt{interp.c} (line 1018).}
\csnippet{MODINT}

\paragraph{Rocq \texttt{Interp.v} (line 707).}
\rocqsnippet{MODINT}

\paragraph{Analysis.}
Same pattern as \texttt{DIVINT}.  Uses \texttt{Z.rem} which matches C's \texttt{\%}
(remainder with truncation toward zero).

\paragraph{Status.} \statusok

\paragraph{Risk.} \risklow


% ─── ANDINT ───────────────────────────────────────────────────────────────────
\instruction{ANDINT}{andint}

\paragraph{PDF Spec.}
\texttt{ANDINT}: Bitwise AND.

\paragraph{OCaml \texttt{interp.c} (line 1024).}
\csnippet{ANDINT}

\paragraph{Rocq \texttt{Interp.v} (line 715).}
\rocqsnippet{ANDINT}

\paragraph{Analysis.}
C: \texttt{(intnat)accu \& (intnat)*sp++}.  Works directly on tagged values
because AND preserves the tag bit.
Rocq: \texttt{Z.land a b}.  Exact match at the value level.

\paragraph{Status.} \statusok

\paragraph{Risk.} \risklow


% ─── ORINT ────────────────────────────────────────────────────────────────────
\instruction{ORINT}{orint}

\paragraph{PDF Spec.}
\texttt{ORINT}: Bitwise OR.

\paragraph{OCaml \texttt{interp.c} (line 1026).}
\csnippet{ORINT}

\paragraph{Rocq \texttt{Interp.v} (line 721).}
\rocqsnippet{ORINT}

\paragraph{Analysis.}
C: \texttt{(intnat)accu | (intnat)*sp++}.  Preserves tag bit.
Rocq: \texttt{Z.lor a b}.  Exact match.

\paragraph{Status.} \statusok

\paragraph{Risk.} \risklow


% ─── XORINT ───────────────────────────────────────────────────────────────────
\instruction{XORINT}{xorint}

\paragraph{PDF Spec.}
\texttt{XORINT}: Bitwise XOR.

\paragraph{OCaml \texttt{interp.c} (line 1028).}
\csnippet{XORINT}

\paragraph{Rocq \texttt{Interp.v} (line 727).}
\rocqsnippet{XORINT}

\paragraph{Analysis.}
C: \texttt{((intnat)accu \^{} (intnat)*sp++) | 1} --- the \texttt{| 1} restores the tag bit.
Rocq: \texttt{Z.lxor a b}.  Exact match (tag bit is not part of our value representation).

\paragraph{Status.} \statusok

\paragraph{Risk.} \risklow


% ─── LSLINT ───────────────────────────────────────────────────────────────────
\instruction{LSLINT}{lslint}

\paragraph{PDF Spec.}
\texttt{LSLINT}: Logical shift left.

\paragraph{OCaml \texttt{interp.c} (line 1030).}
\csnippet{LSLINT}

\paragraph{Rocq \texttt{Interp.v} (line 733).}
\rocqsnippet{LSLINT}

\paragraph{Analysis.}
C: \texttt{((intnat)accu - 1) << Long\_val(*sp++) + 1}.
Rocq: \texttt{Z.shiftl a b}.
The C version works on tagged values: subtract tag, shift, restore tag.
Our version operates on untagged values directly.  Exact match.

\paragraph{Status.} \statusok

\paragraph{Risk.} \risklow


% ─── LSRINT ───────────────────────────────────────────────────────────────────
\instruction{LSRINT}{lsrint}

\paragraph{PDF Spec.}
\texttt{LSRINT}: Logical (unsigned) shift right.

\paragraph{OCaml \texttt{interp.c} (line 1032).}
\csnippet{LSRINT}

\paragraph{Rocq \texttt{Interp.v} (line 739).}
\rocqsnippet{LSRINT}

\paragraph{Analysis.}
C: \texttt{((uintnat)accu) >> Long\_val(*sp++) | 1} --- unsigned shift, restore tag.
Rocq: \texttt{z\_lsr a b} where \texttt{z\_lsr} masks to 63 bits then does
\texttt{Z.shiftr}.  The \texttt{z\_unsigned} function ensures the shift is
logical (fills with zeros from the left).  Exact match.

\paragraph{Status.} \statusok

\paragraph{Risk.} \risklow


% ─── ASRINT ───────────────────────────────────────────────────────────────────
\instruction{ASRINT}{asrint}

\paragraph{PDF Spec.}
\texttt{ASRINT}: Arithmetic (signed) shift right.

\paragraph{OCaml \texttt{interp.c} (line 1034).}
\csnippet{ASRINT}

\paragraph{Rocq \texttt{Interp.v} (line 745).}
\rocqsnippet{ASRINT}

\paragraph{Analysis.}
C: \texttt{((intnat)accu) >> Long\_val(*sp++) | 1} --- signed shift (implementation-defined
but effectively arithmetic on all OCaml platforms), restore tag.
Rocq: \texttt{Z.shiftr a b} --- arithmetic right shift on \texttt{Z}.  Exact match.

\paragraph{Status.} \statusok

\paragraph{Risk.} \risklow


%═══════════════════════════════════════════════════════════════════════════════
\section{Comparisons}
\label{sec:compare}
%═══════════════════════════════════════════════════════════════════════════════

% ─── EQ ───────────────────────────────────────────────────────────────────────
\instruction{EQ}{eq}

\paragraph{PDF Spec.}
\texttt{EQ}: \texttt{accu = (accu == sp[0]); sp++}.  Physical equality.

\paragraph{OCaml \texttt{interp.c} (line 1041).}
\csnippet{EQ}

\paragraph{Rocq \texttt{Interp.v} (line 751).}
\rocqsnippet{EQ}

\paragraph{Analysis.}
C: physical equality (pointer comparison).
Rocq: \texttt{value\_eqb accu sp[0]} which does structural comparison.
For integers, these are identical.  For heap pointers, structural equality
is stricter than physical equality (two equal blocks at different addresses
compare equal structurally but not physically).

This is a known simplification: our structural equality is a conservative
approximation that may return \texttt{true} in cases where physical equality
returns \texttt{false}.  For well-typed OCaml programs, this difference is
only observable with mutable values (which should use \texttt{==} sparingly).

\paragraph{Status.} \statusneedstest{Structural vs. physical equality difference for heap values.}

\paragraph{Risk.} \riskmed --- Correct for integers (the common case) but differs for heap-allocated blocks.


% ─── NEQ ──────────────────────────────────────────────────────────────────────
\instruction{NEQ}{neq}

\paragraph{PDF Spec.}
\texttt{NEQ}: \texttt{accu = (accu != sp[0]); sp++}.

\paragraph{OCaml \texttt{interp.c} (line 1042).}
\csnippet{NEQ}

\paragraph{Rocq \texttt{Interp.v} (line 757).}
\rocqsnippet{NEQ}

\paragraph{Analysis.}
Negation of \texttt{EQ}.  Same caveats.

\paragraph{Status.} \statusneedstest{Same structural vs.\ physical equality issue.}

\paragraph{Risk.} \riskmed


% ─── LTINT ────────────────────────────────────────────────────────────────────
\instruction{LTINT}{ltint}

\paragraph{PDF Spec.}
\texttt{LTINT}: \texttt{accu = (accu < sp[0]); sp++}.

\paragraph{OCaml \texttt{interp.c} (line 1043).}
\csnippet{LTINT}

\paragraph{Rocq \texttt{Interp.v} (line 763).}
\rocqsnippet{LTINT}

\paragraph{Analysis.}
C: \texttt{(intnat)accu < (intnat)*sp++}.  Works directly on tagged values
because the tagging preserves order for signed integers.
Rocq: \texttt{Z.ltb a b}.  Exact match.

\paragraph{Status.} \statusok

\paragraph{Risk.} \risklow


% ─── LEINT ────────────────────────────────────────────────────────────────────
\instruction{LEINT}{leint}

\paragraph{PDF Spec.}
\texttt{LEINT}: \texttt{accu = (accu <= sp[0]); sp++}.

\paragraph{OCaml \texttt{interp.c} (line 1044).}
\csnippet{LEINT}

\paragraph{Rocq \texttt{Interp.v} (line 769).}
\rocqsnippet{LEINT}

\paragraph{Analysis.}
Same pattern.  \texttt{Z.leb}.  Exact match.

\paragraph{Status.} \statusok

\paragraph{Risk.} \risklow


% ─── GTINT ────────────────────────────────────────────────────────────────────
\instruction{GTINT}{gtint}

\paragraph{PDF Spec.}
\texttt{GTINT}: \texttt{accu = (accu > sp[0]); sp++}.

\paragraph{OCaml \texttt{interp.c} (line 1045).}
\csnippet{GTINT}

\paragraph{Rocq \texttt{Interp.v} (line 775).}
\rocqsnippet{GTINT}

\paragraph{Analysis.}
\texttt{Z.gtb}.  Exact match.

\paragraph{Status.} \statusok

\paragraph{Risk.} \risklow


% ─── GEINT ────────────────────────────────────────────────────────────────────
\instruction{GEINT}{geint}

\paragraph{PDF Spec.}
\texttt{GEINT}: \texttt{accu = (accu >= sp[0]); sp++}.

\paragraph{OCaml \texttt{interp.c} (line 1046).}
\csnippet{GEINT}

\paragraph{Rocq \texttt{Interp.v} (line 781).}
\rocqsnippet{GEINT}

\paragraph{Analysis.}
\texttt{Z.geb}.  Exact match.

\paragraph{Status.} \statusok

\paragraph{Risk.} \risklow


% ─── OFFSETINT ────────────────────────────────────────────────────────────────
\instruction{OFFSETINT}{offsetint}

\paragraph{PDF Spec.}
\texttt{OFFSETINT $n$}: \texttt{accu += $n$}.

\paragraph{OCaml \texttt{interp.c} (line 1067).}
\csnippet{OFFSETINT}

\paragraph{Rocq \texttt{Interp.v} (line 787).}
\rocqsnippet{OFFSETINT}

\paragraph{Analysis.}
C: \texttt{accu += *pc << 1} (shift accounts for tagging).
Rocq: \texttt{Val\_int(a + n)}.  Exact match.

\paragraph{Status.} \statusok

\paragraph{Risk.} \risklow


% ─── OFFSETREF ────────────────────────────────────────────────────────────────
\instruction{OFFSETREF}{offsetref}

\paragraph{PDF Spec.}
\texttt{OFFSETREF $n$}: \texttt{Field(accu, 0) += $n$; accu = ()}.
Increments/decrements the value stored in a reference.

\paragraph{OCaml \texttt{interp.c} (line 1071).}
\csnippet{OFFSETREF}

\paragraph{Rocq \texttt{Interp.v} (line 793).}
\rocqsnippet{OFFSETREF}

\paragraph{Analysis.}
C: \texttt{Field(accu, 0) += *pc << 1; accu = Val\_unit}.
Rocq: reads heap, adds $n$ to first field, writes back via \texttt{heap\_update}.
Exact match.

\paragraph{Status.} \statusok

\paragraph{Risk.} \risklow


% ─── ISINT ────────────────────────────────────────────────────────────────────
\instruction{ISINT}{isint}

\paragraph{PDF Spec.}
\texttt{ISINT}: \texttt{accu = Val\_long(accu \& 1)}.  Tests if accu is an integer
(tag bit set).

\paragraph{OCaml \texttt{interp.c} (line 1076).}
\csnippet{ISINT}

\paragraph{Rocq \texttt{Interp.v} (line 805).}
\rocqsnippet{ISINT}

\paragraph{Analysis.}
C: tests tag bit.  Rocq: pattern matches on \texttt{Val\_int}.  Exact match.

\paragraph{Status.} \statusok

\paragraph{Risk.} \risklow


%═══════════════════════════════════════════════════════════════════════════════
\section{Unsigned Comparisons}
\label{sec:unsigned}
%═══════════════════════════════════════════════════════════════════════════════

% ─── ULTINT ───────────────────────────────────────────────────────────────────
\instruction{ULTINT}{ultint}

\paragraph{PDF Spec.}
\texttt{ULTINT}: Unsigned less-than: \texttt{accu = ((uintnat)accu < (uintnat)sp[0])}.

\paragraph{OCaml \texttt{interp.c} (line 1047).}
\csnippet{ULTINT}

\paragraph{Rocq \texttt{Interp.v} (line 942).}
\rocqsnippet{ULTINT}

\paragraph{Analysis.}
C uses \texttt{uintnat} cast for unsigned comparison on tagged values.
Rocq uses \texttt{z\_flip\_sign} which XORs with $2^{62}$ (the sign bit) to
convert signed to unsigned order.  This is a standard trick: flipping the
sign bit maps the signed range $[-2^{62}, 2^{62}-1]$ to $[0, 2^{63}-1]$
monotonically.

The C code compares \emph{tagged} values as unsigned, while our code compares
\emph{untagged} values with sign-flip.  These are equivalent because the tag
bit (lowest bit, always 1) doesn't affect the ordering of tagged integers
when both operands have the tag set.

\paragraph{Status.} \statusok

\paragraph{Risk.} \risklow


% ─── UGEINT ───────────────────────────────────────────────────────────────────
\instruction{UGEINT}{ugeint}

\paragraph{PDF Spec.}
\texttt{UGEINT}: Unsigned greater-or-equal.

\paragraph{OCaml \texttt{interp.c} (line 1048).}
\csnippet{UGEINT}

\paragraph{Rocq \texttt{Interp.v} (line 948).}
\rocqsnippet{UGEINT}

\paragraph{Analysis.}
Same technique as \texttt{ULTINT} with \texttt{Z.geb}.  Exact match.

\paragraph{Status.} \statusok

\paragraph{Risk.} \risklow


%═══════════════════════════════════════════════════════════════════════════════
\section{Branch Comparisons}
\label{sec:bcompare}
%═══════════════════════════════════════════════════════════════════════════════

% ─── BEQ ──────────────────────────────────────────────────────────────────────
\instruction{BEQ}{beq}

\paragraph{PDF Spec.}
\texttt{BEQ $n$ $ofs$}: If $n$ = \texttt{Long\_val(accu)}, branch by $ofs$.

\paragraph{OCaml \texttt{interp.c} (line 1058).}
\csnippet{BEQ}

\paragraph{Rocq \texttt{Interp.v} (line 895).}
\rocqsnippet{BEQ}

\paragraph{Analysis.}
C: \texttt{if (*pc++ == Long\_val(accu)) pc += *pc; else pc++}.
Rocq: matches \texttt{Val\_int a}, checks \texttt{Z.eqb a n}.
Non-integer accu falls through (no branch), which is correct since an
integer literal can never equal a block.

\paragraph{Status.} \statusok

\paragraph{Risk.} \risklow


% ─── BNEQ ─────────────────────────────────────────────────────────────────────
\instruction{BNEQ}{bneq}

\paragraph{PDF Spec.}
\texttt{BNEQ $n$ $ofs$}: Branch if $n \neq$ accu.

\paragraph{OCaml \texttt{interp.c} (line 1059).}
\csnippet{BNEQ}

\paragraph{Rocq \texttt{Interp.v} (line 903).}
\rocqsnippet{BNEQ}

\paragraph{Analysis.}
Negation of \texttt{BEQ}.  Non-integer accu always branches (correct: a block
is never equal to an integer).

\paragraph{Status.} \statusok

\paragraph{Risk.} \risklow


% ─── BLTINT ───────────────────────────────────────────────────────────────────
\instruction{BLTINT}{bltint}

\paragraph{PDF Spec.}
\texttt{BLTINT $n$ $ofs$}: Branch if $n <$ accu (the operand is on the left).

\paragraph{OCaml \texttt{interp.c} (line 1060).}
\csnippet{BLTINT}

\paragraph{Rocq \texttt{Interp.v} (line 914).}
\rocqsnippet{BLTINT}

\paragraph{Analysis.}
C macro: \texttt{if (*pc++ < (intnat)Long\_val(accu))} --- operand $<$ accu.
Rocq: \texttt{Z.ltb n a} --- $n < a$.  Exact match (note: $n$ is the operand,
$a$ is accu).

\paragraph{Status.} \statusok

\paragraph{Risk.} \risklow


% ─── BLEINT ───────────────────────────────────────────────────────────────────
\instruction{BLEINT}{bleint}

\paragraph{PDF Spec.}
\texttt{BLEINT $n$ $ofs$}: Branch if $n \leq$ accu.

\paragraph{OCaml \texttt{interp.c} (line 1061).}
\csnippet{BLEINT}

\paragraph{Rocq \texttt{Interp.v} (line 921).}
\rocqsnippet{BLEINT}

\paragraph{Analysis.}
\texttt{Z.leb n a}.  Exact match.

\paragraph{Status.} \statusok

\paragraph{Risk.} \risklow


% ─── BGTINT ───────────────────────────────────────────────────────────────────
\instruction{BGTINT}{bgtint}

\paragraph{PDF Spec.}
\texttt{BGTINT $n$ $ofs$}: Branch if $n >$ accu.

\paragraph{OCaml \texttt{interp.c} (line 1062).}
\csnippet{BGTINT}

\paragraph{Rocq \texttt{Interp.v} (line 928).}
\rocqsnippet{BGTINT}

\paragraph{Analysis.}
\texttt{Z.gtb n a}.  Exact match.

\paragraph{Status.} \statusok

\paragraph{Risk.} \risklow


% ─── BGEINT ───────────────────────────────────────────────────────────────────
\instruction{BGEINT}{bgeint}

\paragraph{PDF Spec.}
\texttt{BGEINT $n$ $ofs$}: Branch if $n \geq$ accu.

\paragraph{OCaml \texttt{interp.c} (line 1063).}
\csnippet{BGEINT}

\paragraph{Rocq \texttt{Interp.v} (line 935).}
\rocqsnippet{BGEINT}

\paragraph{Analysis.}
\texttt{Z.geb n a}.  Exact match.

\paragraph{Status.} \statusok

\paragraph{Risk.} \risklow


% ─── BULTINT ──────────────────────────────────────────────────────────────────
\instruction{BULTINT}{bultint}

\paragraph{PDF Spec.}
\texttt{BULTINT $n$ $ofs$}: Unsigned branch if $n <$ accu.

\paragraph{OCaml \texttt{interp.c} (line 1064).}
\csnippet{BULTINT}

\paragraph{Rocq \texttt{Interp.v} (line 957).}
\rocqsnippet{BULTINT}

\paragraph{Analysis.}
C: \texttt{(uintnat)*pc++ < (uintnat)Long\_val(accu)}.
Rocq: \texttt{Z.ltb (z\_flip\_sign n) (z\_flip\_sign a)}.
Same \texttt{z\_flip\_sign} technique as \texttt{ULTINT}.  Exact match.

\paragraph{Status.} \statusok

\paragraph{Risk.} \risklow


% ─── BUGEINT ──────────────────────────────────────────────────────────────────
\instruction{BUGEINT}{bugeint}

\paragraph{PDF Spec.}
\texttt{BUGEINT $n$ $ofs$}: Unsigned branch if $n \geq$ accu.

\paragraph{OCaml \texttt{interp.c} (line 1065).}
\csnippet{BUGEINT}

\paragraph{Rocq \texttt{Interp.v} (line 964).}
\rocqsnippet{BUGEINT}

\paragraph{Analysis.}
\texttt{Z.geb (z\_flip\_sign n) (z\_flip\_sign a)}.  Exact match.

\paragraph{Status.} \statusok

\paragraph{Risk.} \risklow


%═══════════════════════════════════════════════════════════════════════════════
\section{Object-Oriented Operations}
\label{sec:oo}
%═══════════════════════════════════════════════════════════════════════════════

% ─── GETMETHOD ────────────────────────────────────────────────────────────────
\instruction{GETMETHOD}{getmethod}

\paragraph{PDF Spec.}
\texttt{GETMETHOD}: \texttt{accu = Field(Field(sp[0], 0), Int\_val(accu))}.
Looks up a method in the object's class table.

\paragraph{OCaml \texttt{interp.c} (line 1084).}
\csnippet{GETMETHOD}

\paragraph{Rocq \texttt{Interp.v} (line 810).}
\rocqsnippet{GETMETHOD}

\paragraph{Analysis.}
C: \texttt{Lookup(sp[0], accu)} = \texttt{Field(Field(sp[0], 0), Int\_val(accu))}.
Rocq: Gets \texttt{field\_or\_heap obj 0} (class table), then
\texttt{field\_or\_heap class\_tbl (Z.to\_nat n)}.  Exact match.

\paragraph{Status.} \statusok

\paragraph{Risk.} \risklow


% ─── GETPUBMET ────────────────────────────────────────────────────────────────
\instruction{GETPUBMET}{getpubmet}

\paragraph{PDF Spec.}
\texttt{GETPUBMET $tag$}: Push object, look up public method by tag.
Uses binary search in the method table (with optional cache).

\paragraph{OCaml \texttt{interp.c} (line 1090).}
\csnippet{GETPUBMET}

\paragraph{Rocq \texttt{Interp.v} (line 835).}
\rocqsnippet{GETPUBMET}

\paragraph{Analysis.}
C: Uses a cache (\texttt{*pc \& Field(meths,1)}) for fast lookup, falling
back to binary search: \texttt{li=3, hi=Field(meths,0)}, binary search on
odd indices.

Rocq: Uses linear scan of pairs \texttt{(closure, tag)} starting after the
2-word header.  This is semantically correct but slower.

\textbf{Difference:} Linear scan vs.\ binary search.  Both produce the same
result for well-formed method tables.  The cache miss path in C also does
binary search, and the Rocq linear scan covers the same elements.

The OCaml method table layout has:
\begin{itemize}
  \item \texttt{meths[0]} = count/size
  \item \texttt{meths[1]} = cache mask
  \item Even indices $\geq 2$: closures
  \item Odd indices $\geq 3$: tags
\end{itemize}

Our scan skips 2 header words and processes pairs, which matches the
$(closure, tag)$ layout at even/odd positions.

\paragraph{Status.} \statusneedstest{Linear scan correctness for OO programs.}

\paragraph{Risk.} \riskmed --- Different algorithm but same result; needs testing with OO code.


% ─── GETDYNMET ────────────────────────────────────────────────────────────────
\instruction{GETDYNMET}{getdynmet}

\paragraph{PDF Spec.}
\texttt{GETDYNMET}: Dynamic method lookup.  \texttt{accu} is the method tag,
\texttt{sp[0]} is the object.

\paragraph{OCaml \texttt{interp.c} (line 1132).}
\csnippet{GETDYNMET}

\paragraph{Rocq \texttt{Interp.v} (line 863).}
\rocqsnippet{GETDYNMET}

\paragraph{Analysis.}
Same as \texttt{GETPUBMET} but without cache.  C uses binary search,
Rocq uses linear scan.  Same correctness argument.

\paragraph{Status.} \statusneedstest{Same as GETPUBMET.}

\paragraph{Risk.} \riskmed


%═══════════════════════════════════════════════════════════════════════════════
\section{Miscellaneous}
\label{sec:misc}
%═══════════════════════════════════════════════════════════════════════════════

% ─── STOP ─────────────────────────────────────────────────────────────────────
\instruction{STOP}{stop}

\paragraph{PDF Spec.}
\texttt{STOP}: Terminate execution, return accu.

\paragraph{OCaml \texttt{interp.c} (line 1147).}
\csnippet{STOP}

\paragraph{Rocq \texttt{Interp.v} (line 971).}
\rocqsnippet{STOP}

\paragraph{Analysis.}
C: restores initial state, returns \texttt{accu}.
Rocq: returns \texttt{Halt accu}.  Exact match.

\paragraph{Status.} \statusok

\paragraph{Risk.} \risklow


% ─── EVENT ────────────────────────────────────────────────────────────────────
\instruction{EVENT}{event}

\paragraph{PDF Spec.}
\texttt{EVENT}: Debugger event hook.

\paragraph{OCaml \texttt{interp.c} (line 1153).}
\csnippet{EVENT}

\paragraph{Rocq \texttt{Interp.v} (line 972).}
\rocqsnippet{EVENT}

\paragraph{Analysis.}
Modeled as no-op (advance PC).  In C, triggers debugger if event count reaches
zero.  Safe to skip since we don't model the debugger.

\paragraph{Status.} \statusok

\paragraph{Risk.} \risklow


% ─── BREAK ────────────────────────────────────────────────────────────────────
\instruction{BREAK}{break}

\paragraph{PDF Spec.}
\texttt{BREAK}: Debugger breakpoint.

\paragraph{Rocq \texttt{Interp.v} (line 973).}
\rocqsnippet{BREAK}

\paragraph{Analysis.}
Modeled as no-op, same as \texttt{EVENT}.

\paragraph{Status.} \statusok

\paragraph{Risk.} \risklow


% ─── PERFORM/RESUME/RESUMETERM/REPERFORMTERM ─────────────────────────────────
\instruction{PERFORM}{perform}

\paragraph{PDF Spec.}
OCaml 5.x effect handlers.  Not present in OCaml 4.14.

\paragraph{Rocq \texttt{Interp.v} (line 975).}
\rocqsnippet{PERFORM}

\paragraph{Analysis.}
Returns \texttt{Error "PERFORM: effects not supported"}.  Correct for OCaml 4.14
(these opcodes don't exist).  If targeting OCaml 5.x, these would need
implementation.

\paragraph{Status.} \statusstub

\paragraph{Risk.} \risklow --- Not used by OCaml 4.14 bytecode.


\instruction{RESUME}{resume}

\paragraph{Analysis.} Same as \texttt{PERFORM}.  Returns error.

\paragraph{Status.} \statusstub

\paragraph{Risk.} \risklow


\instruction{RESUMETERM}{resumeterm}

\paragraph{Analysis.} Same as \texttt{PERFORM}.  Returns error.

\paragraph{Status.} \statusstub

\paragraph{Risk.} \risklow


\instruction{REPERFORMTERM}{reperformterm}

\paragraph{Analysis.} Same as \texttt{PERFORM}.  Returns error.

\paragraph{Status.} \statusstub

\paragraph{Risk.} \risklow


%═══════════════════════════════════════════════════════════════════════════════
\section{Findings Summary}
\label{sec:summary}
%═══════════════════════════════════════════════════════════════════════════════

\subsection{Statistics}

\begin{center}
\begin{tabular}{lc}
\toprule
\textbf{Category} & \textbf{Count} \\
\midrule
Total instructions audited & 100 \\
\statusok & 82 \\
Needs test & 11 \\
Bug found & 1 \\
Stub (not implemented) & 4 \\
Effect handlers (N/A for 4.14) & 4 \\
\bottomrule
\end{tabular}
\end{center}

\subsection{Summary Table}

\begin{longtable}{llll}
\toprule
\textbf{Instruction} & \textbf{Status} & \textbf{Risk} & \textbf{Notes} \\
\midrule
\endhead
\texttt{ACC} & \statusok & \risklow & \\
\texttt{PUSH} & \statusok & \risklow & \\
\texttt{PUSHACC} & \statusok & \risklow & \\
\texttt{POP} & \statusok & \risklow & \\
\texttt{ASSIGN} & \statusok & \risklow & \\
\texttt{ENVACC} & \statusok & \risklow & \\
\texttt{PUSHENVACC} & \statusok & \risklow & \\
\texttt{OFFSETCLOSURE} & \statusok & \risklow & \\
\texttt{PUSHOFFSETCLOSURE} & \statusok & \risklow & \\
\midrule
\texttt{PUSH\_RETADDR} & \statusok & \risklow & \\
\texttt{APPLY} & \statusok & \risklow & \\
\texttt{APPLY1} & \statusok & \risklow & \\
\texttt{APPLY2} & \statusok & \risklow & \\
\texttt{APPLY3} & \statusok & \risklow & \\
\texttt{APPTERM} & \statusok & \risklow & \\
\texttt{APPTERM1} & \statusok & \risklow & \\
\texttt{APPTERM2} & \statusok & \risklow & \\
\texttt{APPTERM3} & \statusok & \risklow & \\
\texttt{RETURN} & \statusok & \risklow & \\
\midrule
\texttt{RESTART} & \statusok & \riskmed & Closure offset arithmetic \\
\texttt{GRAB} & Needs test & \riskhigh{closinfo} & closinfo=0, pc offset \\
\texttt{CLOSURE} & \statusok & \riskmed & closinfo=0 \\
\texttt{CLOSUREREC} & Needs test & \riskhigh{complex} & Infix hdrs, closinfo, push order \\
\midrule
\texttt{GETGLOBAL} & \statusok & \risklow & \\
\texttt{PUSHGETGLOBAL} & \statusok & \risklow & \\
\texttt{GETGLOBALFIELD} & \statusok & \risklow & \\
\texttt{PUSHGETGLOBALFIELD} & \statusok & \risklow & \\
\texttt{SETGLOBAL} & \statusok & \risklow & \\
\midrule
\texttt{ATOM} & \statusok & \risklow & \\
\texttt{PUSHATOM} & \statusok & \risklow & \\
\texttt{MAKEBLOCK} & \statusok & \risklow & \\
\texttt{MAKEBLOCK1} & \statusok & \risklow & \\
\texttt{MAKEBLOCK2} & \statusok & \risklow & \\
\texttt{MAKEBLOCK3} & \statusok & \risklow & \\
\texttt{MAKEFLOATBLOCK} & Needs test & \riskmed & Float model \\
\midrule
\texttt{GETFIELD} & \statusok & \risklow & \\
\texttt{GETFLOATFIELD} & Needs test & \riskmed & Float model \\
\texttt{SETFIELD} & \statusok & \risklow & \\
\texttt{SETFLOATFIELD} & \statusbug{inverted} & \riskhigh{inverted ops} & accu/sp swapped \\
\midrule
\texttt{VECTLENGTH} & \statusok & \risklow & \\
\texttt{GETVECTITEM} & \statusok & \risklow & \\
\texttt{SETVECTITEM} & \statusok & \risklow & \\
\texttt{GETSTRINGCHAR} & \statusok & \risklow & \\
\texttt{GETBYTESCHAR} & \statusok & \risklow & \\
\texttt{SETBYTESCHAR} & \statusok & \risklow & \\
\midrule
\texttt{BRANCH} & \statusok & \risklow & \\
\texttt{BRANCHIF} & \statusok & \risklow & \\
\texttt{BRANCHIFNOT} & \statusok & \risklow & \\
\texttt{SWITCH} & \statusok & \risklow & Depends on loader \\
\texttt{BOOLNOT} & \statusok & \risklow & \\
\midrule
\texttt{PUSHTRAP} & \statusok & \riskmed & Trap link repr differs \\
\texttt{POPTRAP} & \statusok & \risklow & \\
\texttt{RAISE} & \statusok & \riskmed & Trap navigation \\
\texttt{RERAISE} & \statusok & \risklow & \\
\texttt{RAISE\_NOTRACE} & \statusok & \risklow & \\
\texttt{CHECK\_SIGNALS} & \statusok & \risklow & No-op \\
\texttt{C\_CALL} & \statusok & \risklow & Interaction tree \\
\midrule
\texttt{CONSTINT} & \statusok & \risklow & \\
\texttt{PUSHCONSTINT} & \statusok & \risklow & \\
\texttt{NEGINT} & \statusok & \risklow & \\
\texttt{ADDINT} & \statusok & \risklow & \\
\texttt{SUBINT} & \statusok & \risklow & \\
\texttt{MULINT} & \statusok & \risklow & \\
\texttt{DIVINT} & \statusok & \risklow & Raises on /0 \\
\texttt{MODINT} & \statusok & \risklow & Raises on /0 \\
\texttt{ANDINT} & \statusok & \risklow & \\
\texttt{ORINT} & \statusok & \risklow & \\
\texttt{XORINT} & \statusok & \risklow & \\
\texttt{LSLINT} & \statusok & \risklow & \\
\texttt{LSRINT} & \statusok & \risklow & \\
\texttt{ASRINT} & \statusok & \risklow & \\
\midrule
\texttt{EQ} & Needs test & \riskmed & Structural vs physical \\
\texttt{NEQ} & Needs test & \riskmed & Structural vs physical \\
\texttt{LTINT} & \statusok & \risklow & \\
\texttt{LEINT} & \statusok & \risklow & \\
\texttt{GTINT} & \statusok & \risklow & \\
\texttt{GEINT} & \statusok & \risklow & \\
\texttt{OFFSETINT} & \statusok & \risklow & \\
\texttt{OFFSETREF} & \statusok & \risklow & \\
\texttt{ISINT} & \statusok & \risklow & \\
\midrule
\texttt{ULTINT} & \statusok & \risklow & \\
\texttt{UGEINT} & \statusok & \risklow & \\
\texttt{BEQ} & \statusok & \risklow & \\
\texttt{BNEQ} & \statusok & \risklow & \\
\texttt{BLTINT} & \statusok & \risklow & \\
\texttt{BLEINT} & \statusok & \risklow & \\
\texttt{BGTINT} & \statusok & \risklow & \\
\texttt{BGEINT} & \statusok & \risklow & \\
\texttt{BULTINT} & \statusok & \risklow & \\
\texttt{BUGEINT} & \statusok & \risklow & \\
\midrule
\texttt{GETMETHOD} & \statusok & \risklow & \\
\texttt{GETPUBMET} & Needs test & \riskmed & Linear scan vs binary search \\
\texttt{GETDYNMET} & Needs test & \riskmed & Linear scan vs binary search \\
\midrule
\texttt{STOP} & \statusok & \risklow & \\
\texttt{EVENT} & \statusok & \risklow & No-op \\
\texttt{BREAK} & \statusok & \risklow & No-op \\
\texttt{PERFORM} & \statusstub & \risklow & OCaml 5.x only \\
\texttt{RESUME} & \statusstub & \risklow & OCaml 5.x only \\
\texttt{RESUMETERM} & \statusstub & \risklow & OCaml 5.x only \\
\texttt{REPERFORMTERM} & \statusstub & \risklow & OCaml 5.x only \\
\bottomrule
\end{longtable}


%═══════════════════════════════════════════════════════════════════════════════
\section{Known Discrepancies and Action Items}
\label{sec:actions}
%═══════════════════════════════════════════════════════════════════════════════

\subsection{Bug: SETFLOATFIELD operand inversion}
\label{action:setfloatfield}

\textbf{Severity: High.}
In \texttt{interp.c}, \texttt{SETFLOATFIELD} treats \texttt{accu} as the float
array and \texttt{*sp} as the new value.  Our Rocq code swaps these: it treats
the stack top as the array and \texttt{accu} as the new value.

\textbf{Action:} Fix \texttt{Interp.v} to match the C convention.
Specifically, \texttt{SETFLOATFIELD $n$} should:
\begin{enumerate}
  \item Read field $n$ from \texttt{accu} (the float array)
  \item Pop new value from stack
  \item Store new value into field $n$
  \item Set \texttt{accu = Val\_unit}
\end{enumerate}

\textbf{Impact:} Affects programs using mutable float arrays (\texttt{Array.Floatarray}).
Most OCaml programs don't use these directly, but the standard library does
internally.

\subsection{Systemic: closinfo field hardcoded to 0}
\label{action:closinfo}

\textbf{Severity: Medium.}
Affects: \texttt{GRAB}, \texttt{CLOSURE}, \texttt{CLOSUREREC}.

The \texttt{closinfo} field in closure blocks is hardcoded to \texttt{Val\_int 0}
instead of the proper \texttt{Make\_closinfo(arity, start\_env)} value.  In
\texttt{interp.c}, this field encodes:
\begin{itemize}
  \item The arity of the function (bits 56--63 on 64-bit)
  \item The start-of-environment offset (bits 0--55)
\end{itemize}

Currently safe because nothing in our interpreter inspects closinfo.  Would
become a bug if we implement:
\begin{itemize}
  \item Marshaling/serialization (reads closinfo to find env offset)
  \item The \texttt{caml\_apply*} C primitives (read arity)
  \item Debugger integration
\end{itemize}

\textbf{Action:} Compute the proper closinfo value.  Low priority until
one of the above features is needed.

\subsection{Simplification: structural vs.\ physical equality}
\label{action:equality}

\textbf{Severity: Medium.}
Affects: \texttt{EQ}, \texttt{NEQ}.

Our \texttt{value\_eqb} performs structural equality while \texttt{interp.c}
performs physical (pointer) equality.  For immutable values this is equivalent,
but for mutable heap-allocated values, physical equality can distinguish
aliased references from structurally equal copies.

\textbf{Action:} For \texttt{Val\_ptr} values, compare addresses only (not
heap contents).  This is a one-line fix in \texttt{value\_eqb}.

\subsection{Simplification: GETPUBMET/GETDYNMET linear scan}
\label{action:oo}

\textbf{Severity: Low.}
Our linear scan produces correct results but is $O(n)$ per method lookup
instead of $O(\log n)$.  Not a correctness issue, only a performance concern
for extracted code running OO-heavy programs.

\textbf{Action:} Implement binary search to match \texttt{interp.c}.
Low priority.

\subsection{Simplification: float representation}
\label{action:float}

\textbf{Severity: Low--Medium.}
Affects: \texttt{MAKEFLOATBLOCK}, \texttt{GETFLOATFIELD}, \texttt{SETFLOATFIELD},
\texttt{VECTLENGTH}.

Our model stores floats as regular values (no flat/boxed distinction).
The C interpreter uses \texttt{Double\_wosize} words per float and has
special flat storage.  This doesn't affect correctness for integer-only
programs since float operations go through C primitives.

\textbf{Action:} Revisit when float support is prioritized.

\subsection{Note: infix headers in CLOSUREREC}
\label{action:infix}

\textbf{Severity: Low.}
Our infix headers are \texttt{Val\_block Infix\_tag []} (empty block)
while the C version uses \texttt{Make\_header(i*3, Infix\_tag, Caml\_white)}
(a proper GC header encoding the offset).  Since we have no GC and track
offsets via \texttt{Val\_closure}, this is harmless.  Would need revision
if we ever model GC scanning.


\end{document}
